\section{AMC 12 12B Fall 2021}
        \begin{enumerate}
        	\item (\textit{AMC 12 12B Fall 2021}) What is the value of $1234+2341+3412+4123?$
$\textbf{(A)}\: 10{,}000\qquad\textbf{(B)} \: 10{,}010\qquad\textbf{(C)} \: 10{,}110\qquad\textbf{(D)} \: 11{,}000\qquad\textbf{(E)} \: 11{,}110$

 \textbf{Answer: }E



\textbf{Solution 1}

We see that $1, 2, 3,$ and $4$ each appear in the ones, tens, hundreds, and thousands digit exactly once. Since $1+2+3+4=10$, we find that the sum is equal to 
 \[10\cdot(1+10+100+1000)=\boxed{\textbf{(E)} \: 11{,}110}.\]Note that it is equally valid to manually add all four numbers together to get the answer.

 



\textbf{Solution 2}

We have 
 \[1234 + 2341 + 3412 + 4123 = 1111 \left( 1 + 2 + 3 + 4 \right) = \boxed{\textbf{(E)} \: 11{,}110}.\]~Steven Chen (www.professorchenedu.com)

 



\textbf{Solution 3}

We see that the units digit must be $0$, since $4+3+2+1$ is $0$. But every digit from there, will be a $1$ since we have that each time afterwards, we must carry the $1$ from the previous sum. The answer choice that satisfies these conditions is $\boxed{\textbf{(E)} \: 11{,}110}$.

 



\textbf{Solution 4 (Brute Force)}

We can simply add the numbers. $1234 + 2341 + 3412 + 4123 = 11110 \implies \boxed{\textbf{(E)}}$.

 Note: Although this would not take terribly long, it is not recommended to do this in a real contest.

 



\textbf{Solution 5 (Primary School But Also Similar to Solution 4)}

We can simply just add each place value. 

 $1000+2000+3000+4000=10,000$
$+200+300+400+100=1,000$
$+30+40+10+20=100$
$+4+1+2+3=10$
$10,000+1,000+100+10=\boxed{\textbf{(E)} \: 11{,}110}$.

 


	\item (\textit{AMC 12 12B Fall 2021}) What is the area of the shaded figure shown below?

                    \begin{figure}[H]
                        \centering
                        \includegraphics[width=0.5\textwidth]{Problem/AMC_12_12B_Fall_2021_2_Problem_Figure_1.png}
                    \end{figure}
                    
$\textbf{(A)}\: 4\qquad\textbf{(B)} \: 6\qquad\textbf{(C)} \: 8\qquad\textbf{(D)} \: 10\qquad\textbf{(E)} \: 12$

 \textbf{Answer: }B



\textbf{Solution 1 (Area Addition)}

The line of symmetry divides the shaded figure into two congruent triangles, each with base $3$ and height $2.$

 Therefore, the area of the shaded figure is 
 \[2\cdot\left(\frac12\cdot3\cdot2\right)=2\cdot3=\boxed{\textbf{(B)} \: 6}.\]~MRENTHUSIASM ~Wilhelm Z

 



\textbf{Solution 2 (Area Subtraction)}

To find the area of the shaded figure, we subtract the area of the smaller triangle (base $4$ and height $2$) from the area of the larger triangle (base $4$ and height $5$): 
 \[\frac12\cdot4\cdot5-\frac12\cdot4\cdot2=10-4=\boxed{\textbf{(B)} \: 6}.\]~MRENTHUSIASM ~Steven Chen (www.professorchenedu.com) 

 



\textbf{Solution 3 (Shoelace Theorem)}

The consecutive vertices of the shaded figure are $(1,0),(3,2),(5,0),$ and $(3,5).$ By the \href{https://artofproblemsolving.com/wiki/index.php?title=Shoelace\_Theorem}{Shoelace Theorem}, the area is 
 \[\frac12\cdot|(1\cdot2+3\cdot0+5\cdot5+3\cdot0)-(0\cdot3+2\cdot5+0\cdot3+5\cdot1)|=\frac12\cdot12=\boxed{\textbf{(B)} \: 6}.\]~Taco12 ~I-AM-DA-KING

 



\textbf{Solution 4 (Pick's Theorem)}

We have $4$ lattice points in the interior and $6$ lattice points on the boundary. By \href{https://artofproblemsolving.com/wiki/index.php?title=Pick\%27s\_Theorem}{Pick's Theorem}, the area of the shaded figure is 
 \[4+\frac{6}{2}-1 = 4+3-1 = \boxed{\textbf{(B)} \: 6}.\]~danprathab

 


	\item (\textit{AMC 12 12B Fall 2021}) At noon on a certain day, Minneapolis is $N$ degrees warmer than St. Louis. At $4{:}00$ the temperature in Minneapolis has fallen by $5$ degrees while the temperature in St. Louis has risen by $3$ degrees, at which time the temperatures in the two cities differ by $2$ degrees. What is the product of all possible values of $N?$
$\textbf{(A)}\: 10\qquad\textbf{(B)} \: 30\qquad\textbf{(C)} \: 60\qquad\textbf{(D)} \: 100\qquad\textbf{(E)} \: 120$

 \textbf{Answer: }C



\textbf{Solution 1 (Two Variables)}

At noon on a certain day, let $M$ and $L$ be the temperatures (in degrees) in Minneapolis and St. Louis, respectively. It follows that $M=L+N.$

 At $4{:}00,$ we get
 \begin{align*} |(M-5)-(L+3)| &= 2 \\ |M-L-8| &= 2 \\ |N-8| &= 2. \end{align*}
We have two cases:

 
            \begin{enumerate} 
        \item If $N-8=2,$ then $N=10.$

        \item If $N-8=-2,$ then $N=6.$

            \end{enumerate}
            Together, the product of all possible values of $N$ is $10\cdot6=\boxed{\textbf{(C)} \: 60}.$




\textbf{Solution 2 (One Variable)}

At noon on a certain day, the difference of temperatures in Minneapolis and St. Louis is $N$ degrees.

 At $4{:}00,$ the difference of temperatures in Minneapolis and St. Louis is $N-8$ degrees.

 It follows that 
 \[|N-8|=2.\]We continue with the casework in Solution 1 to get the answer $\boxed{\textbf{(C)} \: 60}.$




\textbf{Solution 3}

Let the temperature in Minneapolis be $N$. Then St. Louis is $0$ degrees. $N - 5 = 3 \pm 2$, $N = -10, -6$, $-10 \cdot -6 = \boxed{\textbf{(C)} \: 60}$.

 


	\item (\textit{AMC 12 12B Fall 2021}) Let $n=8^{2022}$. Which of the following is equal to $\frac{n}{4}?$
$\textbf{(A)}\: 4^{1010}\qquad\textbf{(B)} \: 2^{2022}\qquad\textbf{(C)} \: 8^{2018}\qquad\textbf{(D)} \: 4^{3031}\qquad\textbf{(E)} \: 4^{3032}$

 \textbf{Answer: }E



\textbf{Solution 1}

We have 
 \[n=8^{2022}=  \left(8^\frac{2}{3}\right)^{3033}=4^{3033}.\]Therefore, 
 \[\frac{n}4=\boxed{\textbf{(E)} \: 4^{3032}}.\]~kingofpineapplz

 



\textbf{Solution 2}

The requested value is 
 \[\frac{8^{2022}}{4} = \frac{2^{6066}}{4} = \frac{2^{6066}}{2^2} = 2^{6064} = \boxed{\textbf{(E)} \: 4^{3032}}.\]~NH14

 



\textbf{Solution 3}

If we rewrite everything in powers of 2, we get:$n = 8^{2022} = (2^{3})^{2022} = 2^{6064} = (4^{\frac{1}{2}})^{6064} = 4^{\frac{6064}{2}} = \boxed{\textbf{(E)} \: 4^{3032}}.$

 - abed\_nadir (youtube.com/@indianmathguy)

 



\textbf{Solution 4}

$n = 8^{2022}$
${8^{2022}}/{4} = 2 * 8^{2021} = 2 * 2^{6063} = 2^{6064} = \boxed{\textbf{(E)} \: 4^{3032}}.$

 -sphericalfox

 


	\item (\textit{AMC 12 12B Fall 2021}) Call a fraction $\frac{a}{b}$, not necessarily in the simplest form, \textit{special} if $a$ and $b$ are positive integers whose sum is $15$. How many distinct integers can be written as the sum of two, not necessarily different, special fractions?


 $\textbf{(A)}\ 9 \qquad\textbf{(B)}\  10 \qquad\textbf{(C)}\  11 \qquad\textbf{(D)}\ 12 \qquad\textbf{(E)}\ 13$

 \textbf{Answer: }C



\textbf{Solution 1}

The special fractions are 
 \[\frac{1}{14},\frac{2}{13},\frac{3}{12},\frac{4}{11},\frac{5}{10},\frac{6}{9},\frac{7}{8},\frac{8}{7},\frac{9}{6},\frac{10}{5},\frac{11}{4},\frac{12}{3},\frac{13}{2},\frac{14}{1}.\]We rewrite them in the simplest form: 
 \[\frac{1}{14},\frac{2}{13},\frac{1}{4},\frac{4}{11},\frac{1}{2},\frac{2}{3},\frac{7}{8},1\frac{1}{7},1\frac{1}{2},2,2\frac{3}{4},4,6\frac{1}{2},14.\]Note that two unlike fractions in the simplest form cannot sum to an integer. So, we only consider the fractions whose denominators appear more than once: 
 \[\frac{1}{4},\frac{1}{2},1\frac{1}{2},2,2\frac{3}{4},4,6\frac{1}{2},14.\]For the set $\{2,4,14\},$ two elements (not necessarily different) can sum to $4,6,8,16,18,28.$

 For the set $\left\{\frac{1}{2},1\frac{1}{2},6\frac{1}{2}\right\},$ two elements (not necessarily different) can sum to $1,2,3,7,8,13.$

 For the set $\left\{\frac{1}{4},2\frac{3}{4}\right\},$ two elements (not necessarily different) can sum to $3.$

 Together, there are $\boxed{\textbf{(C)}\  11}$ distinct integers that can be written as the sum of two, not necessarily different, special fractions: 
 \[1,2,3,4,6,7,8,13,16,18,28.\]~KingRavi ~samrocksnature ~Wilhelm Z ~MRENTHUSIASM ~Steven Chen (www.professorchenedu.com)

 



\textbf{Solution 2}

Let $a=15-b,$ so the special fraction is 
 \[\frac ab = \frac{15-b}{b} = \frac{15}{b}-1.\]We can ignore the $-1$ part and only focus on $\frac{15}{b}.$

 The integers are $\frac{15}{1},\frac{15}{3},\frac{15}{5},$ which are $15,5,3,$ respectively. We get $30,20,18,10,8,6$ from this group of numbers.

 The halves are $\frac{15}{2},\frac{15}{6},\frac{15}{10},$ which are $7\frac12,2\frac12,1\frac12,$ respectively. We get $15,10,9,5,4,3$ from this group of numbers.

 The quarters are $\frac{15}{4},\frac{15}{12},$ which are $3\frac34,1\frac14,$ respectively. We get $5$ from this group of numbers.

 Note that $10$ and $5$ each appear twice. Therefore, the answer is $\boxed{\textbf{(C)}\  11}.$




\textbf{Solution 3 (Casework)}

We split this up into two cases:

 \textbf{Case 1: integer + integer}

 The whole numbers we have are $\frac{10}{5}$ (or $2$), $\frac{12}{3}$ (or $4$), and $\frac{14}{1}$ (or $14$). There are $\dbinom{3}{2}=3$ ways to choose different-numbered pairs and $3$ ways to choose the same-numbered pairs. So, $3+3=6$.

 \textbf{Case 2: fraction + fraction}

 The fractions we have are $\frac{5}{10}$ (or $\frac{1}{2}$), $\frac{9}{6}$ (or $\frac{3}{2}$), and $\frac{13}{2}$. Similarly, there are $\dbinom{3}{2}=3$ ways to choose different-numbered pairs and $3$ ways to choose the same-numbered pairs. So, $3+3=6$.

 Thus, $6+6=12$.

 So now you would just go ahead and innocently choose $\textbf{(D) }12$, right? No! We overcounted $8$, as $\frac{13}{2}+\frac96=\frac{12}{3}+\frac{12}{3}=8$. Therefore, the correct answer is actually $12-1=\boxed{\textbf{(C)}\  11}$.

 


	\item (\textit{AMC 12 12B Fall 2021}) The greatest prime number that is a divisor of $16{,}384$ is $2$ because $16{,}384 = 2^{14}$. What is the sum of the digits of the greatest prime number that is a divisor of $16{,}383$?


 $\textbf{(A)} \: 3\qquad\textbf{(B)} \: 7\qquad\textbf{(C)} \: 10\qquad\textbf{(D)} \: 16\qquad\textbf{(E)} \: 22$

 \textbf{Answer: }C



\textbf{Solution}

We have
 \begin{align*} 16383 & = 2^{14} - 1 \\ & = \left( 2^7 + 1 \right) \left( 2^7 - 1 \right) \\ & = 129 \cdot 127 \\ \end{align*}


 Since $129$ is composite, $127$ is the largest prime which can divide $16383$. The sum of $127$'s digits is $1+2+7=\boxed{\textbf{(C) }10}$.

 \textbf{Note} Note that you can quickly tell that $2^7 -1$ is prime because it is a \href{https://artofproblemsolving.com/wiki/index.php?title=Mersenne\_prime}{Mersenne prime}.

 


	\item (\textit{AMC 12 12B Fall 2021}) Which of the following conditions is sufficient to guarantee that integers $x$, $y$, and $z$ satisfy the equation

 \[x(x-y)+y(y-z)+z(z-x) = 1?\]
$\textbf{(A)} \: x>y$ and $y=z$
$\textbf{(B)} \: x=y-1$ and $y=z-1$
$\textbf{(C)} \: x=z+1$ and $y=x+1$
$\textbf{(D)} \: x=z$ and $y-1=x$
$\textbf{(E)} \: x+y+z=1$

 \textbf{Answer: }D



\textbf{Solution 1 (Completing the Square)}

It is obvious $x$, $y$, and $z$ are symmetrical. We are going to solve the problem by Completing the Square.

 $x ^ 2 + y ^ 2 + z ^ 2 - xy - yz - zx = 1$
$2x ^ 2 + 2y ^ 2 + 2z ^ 2 - 2xy - 2yz - 2zx = 2$
$(x-y)^2 + (y-z)^2 + (z-x)^2 = 2$

 Because 

 


	\item (\textit{AMC 12 12B Fall 2021}) The product of the lengths of the two congruent sides of an obtuse isosceles triangle is equal to the product of the base and twice the triangle's height to the base. What is the measure, in degrees, of the vertex angle of this triangle?


 $\textbf{(A)} \: 105 \qquad\textbf{(B)} \: 120 \qquad\textbf{(C)} \: 135 \qquad\textbf{(D)} \: 150 \qquad\textbf{(E)} \: 165$

 \textbf{Answer: }D



\textbf{Solution 1 (Area)}

Let the lengths of the two congruent sides of the triangle be $x$, then the product desired is $x^2$.

 Notice that the product of the base and twice the height is $4$ times the area of the triangle.

 Set the vertex angle to be $a$, we derive the equation:

 $x^2=4\left(\frac{1}{2}x^2\sin(a)\right)$
$\sin(a)=\frac{1}{2}$

 As the triangle is obtuse, $a=150^\circ$ only. We get $\boxed{\textbf{(D)} \ 150}.$




\textbf{Solution 2}

Denote by $a$ the length of each congruent side. Denote by $\theta$ the degree measure of each acute angle.Denote by $\phi$ the degree measure of the obtuse angle.

 Hence, this problem tells us the following relationship:
 \[ a^2 = 2 a \cos \theta \cdot 2 a \sin \theta . \]

 Hence,
 \begin{align*} 1 & = 2 \cdot 2 \sin \theta \cos \theta \\ & = 2 \sin 2 \theta \\ & = 2 \sin \left( 180^\circ - 2 \theta \right) \\ & = 2 \sin \phi . \end{align*}


 Hence, $\phi = 150^\circ$.

 Therefore, the answer is $\boxed{\textbf{(D) }150}$.

 


 


	\item (\textit{AMC 12 12B Fall 2021}) Triangle $ABC$ is equilateral with side length $6$. Suppose that $O$ is the center of the inscribed
circle of this triangle. What is the area of the circle passing through $A$, $O$, and $C$?


 $\textbf{(A)} \: 9\pi \qquad\textbf{(B)} \: 12\pi \qquad\textbf{(C)} \: 18\pi \qquad\textbf{(D)} \: 24\pi \qquad\textbf{(E)} \: 27\pi$

 \textbf{Answer: }B



\textbf{Solution 1 (Cosine Rule)}

Construct the circle that passes through $A$, $O$, and $C$, centered at $X$. 

 Also notice that $\overline{OA}$ and $\overline{OC}$ are the angle bisectors of angle $\angle BAC$ and $\angle BCA$ respectively. We then deduce $\angle AOC=120^\circ$.

 Consider another point $M$ on Circle $X$ opposite to point $O$. 

 As $AOCM$ is an inscribed quadrilateral of Circle $X$, $\angle AMC=180^\circ-120^\circ=60^\circ$.

 Afterward, deduce that $\angle AXC=2·\angle AMC=120^\circ$.

 By the Cosine Rule, we have the equation: (where $r$ is the radius of circle $X$)

 $2r^2(1-\cos(120^\circ))=6^2$
$r^2=12$

 The area is therefore $\pi r^2 = \boxed{\textbf{(B)}\ 12\pi}$.

 



\textbf{Solution 2}

We have $\angle AOC = 120^\circ$.

 Denote by $R$ the circumradius of $\triangle AOC$.In $\triangle AOC$, the law of sines implies
 \[ 2 R = \frac{AC}{\sin \angle AOC} = 4 \sqrt{3} . \]

 Hence, the area of the circumcircle of $\triangle AOC$ is
 \[ \pi R^2 = 12 \pi . \]

 Therefore, the answer is $\boxed{\textbf{(B) }12 \pi}$.

 



\textbf{Solution 3}

As in the previous solution, construct the circle that passes through $A$, $O$, and $C$, centered at $X$. Let $Y$ be the intersection of $\overline{OX}$ and $\overline{AB}$.

 Note that since $\overline{OA}$ is the angle bisector of $\angle BAC$ that $\angle OAC=30^\circ$. Also by symmetry, $\overline{OX}$ $\perp$ $\overline{AB}$ and $AY = 3$. Thus $\tan(30^\circ) = \frac{OY}{3}$ so $OY = \sqrt{3}$.

 Let $r$ be the radius of circle $X$, and note that $AX = OX = r$. So $\triangle AYX$ is a right triangle with legs of length $3$ and $r - \sqrt{3}$ and hypotenuse $r$. By Pythagoras, $3^2 + (r - \sqrt{3})^2 = r^2$. So $r = 2\sqrt{3}$.

 Thus the area is $\pi r^2 = \boxed{\textbf{(B)}\ 12\pi}$.

 -SharpeMind

 



\textbf{Solution 5 (SIMPLE)}

The semiperimeter is $\frac{6+6+6}{2}=9$ units.The area of the triangle is $9\sqrt{3}$ units squared.  By the formula that says that the area of the triangle is its semiperimeter times its inradius, the inradius $r=\sqrt{3}$. As $\angle{AOC}=120^\circ$, we can form an altitude from point $O$ to side $AC$ at point $M$, forming two 30-60-90 triangles. As $CM=MA=3$, we can solve for $OC=2\sqrt{3}$. Now, call the center of the circle we are looking for $X$. From here we can now utilize the fact that since the inscribed angle $\angle{AOC}=120^\circ$ that the arclength that the angle describes is $240^\circ$. This leaves us $120^\circ$ ($360-240$) left to be formed by the central angle, which by the Central Angle Theorem must be $120^\circ$ as well. Note that now we have a parallelogram with the side $OC$ (which is $2\sqrt{3}$) being opposite and therefor equivalent to the radius $XA$. Now, the area of the circle is just $\pi*(2*\sqrt{3})^2 = 12\pi$. Select $\boxed{B}$.

 \textbf{Note For Solution 5} One must be careful as one can easily accidentally get this problem CORRECT by prematurely saying the radius is $2\sqrt{3}$ BEFORE proving that $OC$ is a part of a parallelogram. Otherwise, your statement is simply saying that the Circumcircle's area is $12\pi$. Meaning that although your answer was correct the reasoning was not.

 



\textbf{Solution 6 (Ptolemy)}

Call the diameter of the circle $d$. If we extend points $A$ and $C$ to meet at a point on the circle and call it $E$, then $\bigtriangleup OAE=\bigtriangleup OCE$ . Note that both triangles are right, since their hypotenuse is the diameter of the circle. Therefore, $CE=AE=\sqrt{d^2-12}$. We know this since $OC=OA=OB$ and $OC$ is the hypotenuse of a $30-60-90$ right triangle, with the longer leg being $\frac{6}{2}=3$ so $OC=2\sqrt{3}$. Applying Ptolemy's Theorem on cyclic quadrilateral $OCEA$, we get $2({\sqrt{d^2-12}})\cdot{2\sqrt{3}}=6d$. Squaring and solving we get $d^2=48 \Longrightarrow (2r)^2=48$ so $r^2=12$. Therefore, the area of the circle is $\boxed{12\pi}$



	\item (\textit{AMC 12 12B Fall 2021}) What is the sum of all possible values of $t$ between $0$ and $360$ such that the triangle in the coordinate plane whose vertices are 
 \[(\cos 40^\circ,\sin 40^\circ), (\cos 60^\circ,\sin 60^\circ), \text{ and } (\cos t^\circ,\sin t^\circ)\]
is isosceles? 


 $\textbf{(A)} \: 100 \qquad\textbf{(B)} \: 150 \qquad\textbf{(C)} \: 330 \qquad\textbf{(D)} \: 360 \qquad\textbf{(E)} \: 380$

 \textbf{Answer: }E



\textbf{Solution}

Let $A = (\cos 40^{\circ}, \sin 40^{\circ}), B = (\cos 60^{\circ}, \sin 60^{\circ}),$ and $C = (\cos t^{\circ}, \sin t^{\circ}).$ We apply casework to the legs of isosceles $\triangle ABC:$

            \begin{enumerate} 
        \item $AB=AC$Note that $A$ must be the midpoint of $\widehat{BC}.$ It follows that $C = (\cos 20^{\circ}, \sin 20^{\circ}),$ so $t=20.$


        \item $BA=BC$Note that $B$ must be the midpoint of $\widehat{AC}.$ It follows that $C = (\cos 80^{\circ}, \sin 80^{\circ}),$ so $t=80.$


        \item $CA=CB$Note that $C$ must be the midpoint of $\widehat{AB}.$ It follows that $C = (\cos 50^{\circ}, \sin 50^{\circ})$ or $C = (\cos 230^{\circ}, \sin 230^{\circ}),$ so $t=50$ or $t=230.$


            \end{enumerate}
            Together, the sum of all such possible values of $t$ is $20+80+50+230=\boxed{\textbf{(E)} \: 380}.$
<u>\textbf{Remark}</u>

 The following diagram shows all possible locations of $C:$

                    \begin{figure}[H]
                        \centering
                        \includegraphics[width=0.5\textwidth]{Solution/AMC_12_12B_Fall_2021_10_Solution_Figure_1.png}
                    \end{figure}
                    



	\item (\textit{AMC 12 12B Fall 2021}) Una rolls $6$ standard $6$-sided dice simultaneously and calculates the product of the $6{ }$ numbers obtained. What is the probability that the product is divisible by $4?$
$\textbf{(A)}\: \frac34\qquad\textbf{(B)} \: \frac{57}{64}\qquad\textbf{(C)} \: \frac{59}{64}\qquad\textbf{(D)} \: \frac{187}{192}\qquad\textbf{(E)} \: \frac{63}{64}$

 \textbf{Answer: }C



\textbf{Solution}

We will use complementary counting to find the probability that the product is not divisible by $4$. Then, we can find the probability that we want by subtracting this from 1. We split this into $2$ cases. 

 Case 1: The product is not divisible by $2$.

 We need every number to be odd, and since the chance we roll an odd number is $\frac12,$ our probability is $\left(\frac12\right)^6=\frac1{64}.$

 Case 2: The product is divisible by $2$, but not by $4$.

 We need $5$ numbers to be odd, and one to be divisible by $2$, but not by $4$. There is a $\frac12$ chance that an odd number is rolled, a $\frac13$ chance that we roll a number satisfying the second condition (only $2$ and $6$ work), and $6$ ways to choose the order in which the even number appears.

 Our probability is $\left(\frac12\right)^5\left(\frac13\right)\cdot6=\frac1{16}.$

 Therefore, the probability the product is not divisible by $4$ is $\frac1{64}+\frac1{16}=\frac{5}{64}$.

 Our answer is $1-\frac{5}{64}=\boxed{\textbf{(C)}\ \frac{59}{64}}$.

 



\textbf{Solution 2(Pretty much same as above)}

Complementary counting implies that we can have either 0 factors of 2 or just 1 factor of 2.

 Case 1: 0 factors of 2

 The only odd numbers are 1, 3, 5 on a die so we have these 3 odd numbers with 6 positions so $3^6$ ways here.

 Case 2: 1 factor of 2

 The only even numbers that have only one factor of 2 on a die are 2, 6. On the 6 positions, there are 6 ways to place one of these even numbers and the rest of the 5 positions must be one of the three odd numbers. Since we could have either 2 or 6, we have

 $2 \cdot 6 \cdot 3^5 = 4 \cdot 3^6$. 

 Thus, we have $3^6 + 4 \cdot 3^6 = 5 \cdot 3^6$ cases that don't work.

 By complementary counting we get

 $1 - \frac{5 \cdot 3^6}{6^6} = 1 - \frac{5}{2^6} = \boxed{\frac{59}{64}}$



	\item (\textit{AMC 12 12B Fall 2021}) For $n$ a positive integer, let $f(n)$ be the quotient obtained when the sum of all positive divisors of $n$ is divided by $n.$ For example, 
 \[f(14)=(1+2+7+14)\div 14=\frac{12}{7}\]
What is $f(768)-f(384)?$
$\textbf{(A)}\ \frac{1}{768} \qquad\textbf{(B)}\ \frac{1}{192} \qquad\textbf{(C)}\ 1 \qquad\textbf{(D)}\ \frac{4}{3} \qquad\textbf{(E)}\ \frac{8}{3}$

 \textbf{Answer: }B



\textbf{Solution 1}

The prime factorizations of $768$ and $384$ are $2^8\cdot3$ and $2^7\cdot3,$ respectively. Note that $f(n)$ is the sum of all fractions of the form $\frac 1d,$ where $d$ is a positive divisor of $n.$ By geometric series, it follows that
 \begin{alignat*}{8} f(768)&=\left(\sum\_{k=0}\textasciicircum{}{8}\frac{1}{2\textasciicircum{}k}\right)+\left(\sum\_{k=0}\textasciicircum{}{8}\frac{1}{2\textasciicircum{}k\cdot3}\right)&&=\frac{511}{256}+\frac{511}{768}&&=\frac{2044}{768}, \\ f(384)&=\left(\sum\_{k=0}\textasciicircum{}{7}\frac{1}{2\textasciicircum{}k}\right)+\left(\sum\_{k=0}\textasciicircum{}{7}\frac{1}{2\textasciicircum{}k\cdot3}\right)&&=\frac{255}{128}+\frac{255}{384}&&=\frac{1020}{384}. \end{alignat*}Therefore, the answer is $f(768)-f(384)=\boxed{\textbf{(B)}\ \frac{1}{192}}.$




\textbf{Solution 2}

The prime factorization of $384$ is $2^7\cdot3,$ so each of its positive divisors is of the form $2^m$ or $2^m\cdot3$ for some integer $m$ such that $0\leq m\leq7.$ We will use this fact to calculate the sum of all its positive divisors. Note that 
 \[2^m + 2^m\cdot3 = 2^m\cdot(1+3) = 2^m\cdot4 = 2^m\cdot2^2 = 2^{m+2}\] is the sum of the two forms of positive divisors for all such $m.$ By geometric series, the sum of all positive divisors of $384$ is 
 \[\sum_{k=2}^{9}2^k = \frac{2^{10}-2^2}{2-1} = 1020,\] from which $f(384) = \frac{1020}{384}.$ Similarly, since the prime factorization of $768$ is $2^8 \cdot 3,$ we have $f(768) = \frac{2044}{768}.$ 

 Therefore, the answer is $f(768)-f(384)=\boxed{\textbf{(B)}\ \frac{1}{192}}.$




\textbf{Solution 3}

Let $\sigma(n)$ denotes the sum of all positive divisors of $n,$ so $f(n)=\sigma(n)\div n.$

 Suppose that $n=\prod_{i=1}^{k}p_i^{e_i}$ is the prime factorization of $n.$ Since $\sigma(n)$ is multiplicative, we have 
 \[\sigma(n)=\sigma\left(\prod_{i=1}^{k}p_i^{e_i}\right)=\prod_{i=1}^{k}\sigma\left(p_i^{e_i}\right)=\prod_{i=1}^{k}\left(\sum_{j=0}^{e_i}p_i^j\right)=\prod_{i=1}^{k}\frac{p_i^{e_i+1}-1}{p_i-1}\] by geometric series.

 The prime factorizations of $768$ and $384$ are $2^8\cdot3$ and $2^7\cdot3,$ respectively. Note that
 \begin{alignat*}{8} f(768) &= \left(\frac{2\textasciicircum{}9-1}{2-1}\cdot\frac{3\textasciicircum{}2-1}{3-1}\right)\div768 &&= \frac{2044}{768}, \\ f(384) &= \left(\frac{2\textasciicircum{}8-1}{2-1}\cdot\frac{3\textasciicircum{}2-1}{3-1}\right)\div384 &&= \frac{1020}{384}. \end{alignat*}Therefore, the answer is $f(768)-f(384)=\boxed{\textbf{(B)}\ \frac{1}{192}}.$




\textbf{Solution 4}

Note that
 \begin{align*} 384 &= 2^7\cdot3, \\ 768 &= 2^8\cdot3. \end{align*}
Both numbers take the form $2^n\cdot3$. Define a function $S(n)$ as being the sum of the factors for any number $2^n\cdot3$, where $n$ is a nonnegative integer.

 If you list out the factors of a number $2^n\cdot3$, all the factors are even except for $1$ and $3$. So to produce a list of factors for the number $2^{n+1}\cdot3$, you can multiply all the factors on the first list by $2$ and then add on $1$ and $3$. Thus, $S(n) = 2S(n-1) + 4$.

 We need to find $\frac{S(8)}{768} - \frac{S(7)}{384}$. This can be written as $\frac{2S(7)+4-2S(7)}{768}$. Our recursive formula states that $S(n) = 2S(n-1) + 4$ so $S(n) - 2S(n-1) = 4$. Thus, the fraction simplifies to $\frac{4}{768} = \boxed{\textbf{(B)}\ \frac{1}{192}}$. Note that we never needed to figure out what $S(8)$ or $S(7)$ were, sparing us a lot of calculation.

 



\textbf{Solution 5}

We know that
 \begin{align*} f(768) &= \frac{1 + 2 + 3 + 4 + \cdots + 768}{768}, \\ f(384) &= \frac{1 + 2 + 3 + 4 + \cdots + 384}{384}. \end{align*}
We want to find the difference of these so
 \[\frac{1 + 2 + 3 + 4 + \cdots + 768}{768} - \frac{ 2(1 + 2 + 3 + 4 + \cdots + 384)}{2\cdot384} = \frac{1 + 2 + 3 + 4 + \cdots + 768}{768} - \frac{ 2 + 4 + 6 + 8 + \cdots + 768}{768}.\]We are only left with the even divisors of $768$ on the fraction on the right so the answer would be the sum of all odd divisors  of $768$ ($1$ and $3$) over $768$:
 \[\frac{1 + 3}{768}=\boxed{\textbf{(B)}\ \frac{1}{192}}.\]~pengf

 



\textbf{Solution 6}

As in Solution 1, the prime factorization of $768$ is $2^8 \cdot 3$, and the prime factorization of $384$ is $2^7 \cdot 3$. Recalling the sum of factors formula and sum of the powers of 2 formula, we find that the sum of the factors of $768$ is $(1+2^1+2^2+2^3+2^4+2^5+2^6+2^7+2^8)(1+3^1)=(2^9-1)(4)=(511)(4)$, and the sum of factors of $384$ is $1+2^1+2^2+2^3+2^4+2^5+2^6+2^7)(1+3^1)=(2^8-1)(4)=(255)(4)$. Then $f(768)=\frac{(511)(4)}{768}=\frac{511}{192}$, and $f(384)=\frac{(255)(4)}{384}=\frac{510}{192}$. So, $f(768)-f(384)=\frac{511}{192}-\frac{510}{192}=\boxed{\textbf{(B)}\ \frac{1}{192}}$.

 \textbf{Note} Note that you can cut away some of the arithmetic by realizing that $(2^8-1) \cdot 2$ is 1 less than $2^9-1$.~bird987

 



\textbf{Solution 7}

Observe that $f(n)$ equals the sum of the reciprocals of the positive divisors of $n$. Since 768 is a multiple of 384, every divisor of 384 will also be a divisor of 768. Then $f(768)-f(384)$ equals the sum of the reciprocals of the positive divisors of 768 which are not divisors of 384. Since $768 = 2^8 \cdot 3$ and $384 = 2^7 \cdot 3$, the divisors of 768 which are not divisors of 384 are $2^8 = 256$ and $2^8 \cdot 3 = 768$. Therefore

 


	\item (\textit{AMC 12 12B Fall 2021}) Let $c = \frac{2\pi}{11}.$ What is the value of

 \[\frac{\sin 3c \cdot \sin 6c \cdot \sin 9c \cdot \sin 12c \cdot \sin 15c}{\sin c \cdot \sin 2c \cdot \sin 3c \cdot \sin 4c \cdot \sin 5c}?\]
$\textbf{(A)}\ {-}1 \qquad\textbf{(B)}\ {-}\frac{\sqrt{11}}{5} \qquad\textbf{(C)}\ \frac{\sqrt{11}}{5} \qquad\textbf{(D)}\ \frac{10}{11} \qquad\textbf{(E)}\ 1$

 \textbf{Answer: }E



\textbf{Solution}

Plugging in $c$, we get
 \[\frac{\sin 3c \cdot \sin 6c \cdot \sin 9c \cdot \sin 12c \cdot \sin 15c}{\sin c \cdot \sin 2c \cdot \sin 3c \cdot \sin 4c \cdot \sin 5c}=\frac{\sin \frac{6\pi}{11} \cdot \sin \frac{12\pi}{11} \cdot \sin \frac{18\pi}{11} \cdot \sin \frac{24\pi}{11} \cdot \sin \frac{30\pi}{11}}{\sin \frac{2\pi}{11} \cdot \sin \frac{4\pi}{11} \cdot \sin \frac{6\pi}{11} \cdot \sin \frac{8\pi}{11} \cdot \sin \frac{10\pi}{11}}.\]Since $\sin(x+2\pi)=\sin(x),$ $\sin(2\pi-x)=\sin(-x),$ and $\sin(-x)=-\sin(x),$ we get
 \[\frac{\sin \frac{6\pi}{11} \cdot \sin \frac{12\pi}{11} \cdot \sin \frac{18\pi}{11} \cdot \sin \frac{24\pi}{11} \cdot \sin \frac{30\pi}{11}}{\sin \frac{2\pi}{11} \cdot \sin \frac{4\pi}{11} \cdot \sin \frac{6\pi}{11} \cdot \sin \frac{8\pi}{11} \cdot \sin \frac{10\pi}{11}}=\frac{\sin \frac{6\pi}{11} \cdot \sin \frac{-10\pi}{11} \cdot \sin \frac{-4\pi}{11} \cdot \sin \frac{2\pi}{11} \cdot \sin \frac{8\pi}{11}}{\sin \frac{2\pi}{11} \cdot \sin \frac{4\pi}{11} \cdot \sin \frac{6\pi}{11} \cdot \sin \frac{8\pi}{11} \cdot \sin \frac{10\pi}{11}}=\boxed{\textbf{(E)}\ 1}.\]




\textbf{Solution 2}

Eisenstein used such a quotient in his proof of \href{https://artofproblemsolving.com/wiki/index.php?title=Quadratic\_reciprocity}{quadratic reciprocity}. Let $c=\frac{2\pi}{p}$ where $p$ is an odd prime number and $q$ is any integer.

 Then $\dfrac{\sin(qc)\sin(2qc)\cdots\sin(\frac{p-1}{2}qc)}{\sin(c)\sin(2c)\cdots\sin(\frac{p-1}{2}c)}$ is the Legendre symbol $\left(\frac{q}{p}\right)$. Legendre symbol is calculated using quadratic reciprocity which is $\left(\frac{p}{q}\right)\left(\frac{q}{p}\right)=(-1)^{\frac{p-1}{2}\frac{q-1}{2}}$. The Legendre symbol $\left(\frac{3}{11}\right)=(-1)\left(\frac{11}{3}\right)=(-1)\left(\frac{-1}{3}\right)=(-1)(-1)=\boxed{\textbf{(E)}\ 1}$




\textbf{Solution 3}

We have that $5^2 \equiv 3 \pmod{11}$, so 3 is a quadratic residue mod 11. For quadratic residues, their Legendre symbol which we know is the answer from Solution 2 is $\boxed{\textbf{(E)}\ 1}$




\textbf{Solution 4}

We have $\zeta$ be an 11th primitive root of unity. Then the quotient becomes 
 \[\frac{(\zeta^3-\zeta^{-3})(\zeta^6-\zeta^{-6})(\zeta^9-\zeta^{-9})(\zeta^{12}-\zeta^{-12})(\zeta^{15}-\zeta^{-15})}{(\zeta^1-\zeta^{-1})(\zeta^2-\zeta^{-2})(\zeta^3-\zeta^{-3})(\zeta^{4}-\zeta^{-4})(\zeta^{5}-\zeta^{-5})}\] which we can use modular arithmetic to become 
 \[\frac{(\zeta^3-\zeta^{-3})(\zeta^{-5}-\zeta^{5})(\zeta^{-2}-\zeta^{2})(\zeta^{1}-\zeta^{-1})(\zeta^{4}-\zeta^{-4})}{(\zeta^1-\zeta^{-1})(\zeta^2-\zeta^{-2})(\zeta^3-\zeta^{-3})(\zeta^{4}-\zeta^{-4})(\zeta^{5}-\zeta^{-5})}\] and we see that is $\boxed{\textbf{(E)}\ 1}$




\textbf{Solution 5 (guess)}

Since we can't calculate the exact value of $\sin{\frac{2\pi}{11}},$ it's unlikely the answer will have radicals in it. This leaves us with $-1, 1,$ and $\frac{10}{11}.$ The fraction is too precise for it to be a likely answer, so that leaves us with $1$ or $-1.$ Note that if the angle is greater than $\pi,$ the sine will be negative, and if it's less than $\pi$ it'll be positive. The numerator and denominator are therefore both positive, so we eliminate $-1$ to get $\boxed{\textbf{(E)}\ 1}.$



	\item (\textit{AMC 12 12B Fall 2021}) Suppose that $P(z), Q(z)$, and $R(z)$ are polynomials with real coefficients, having degrees $2$, $3$, and $6$, respectively, and constant terms $1$, $2$, and $3$, respectively. Let $N$ be the number of distinct complex numbers $z$ that satisfy the equation $P(z) \cdot Q(z)=R(z)$. What is the minimum possible value of $N$?


 $\textbf{(A)}\: 0\qquad\textbf{(B)} \: 1\qquad\textbf{(C)} \: 2\qquad\textbf{(D)} \: 3\qquad\textbf{(E)} \: 5$

 \textbf{Answer: }B



\textbf{Solution}

The answer cannot be $0,$ as every nonconstant polynomial has at least $1$ distinct complex root (Fundamental Theorem of Algebra). Since $P(z) \cdot Q(z)$ has degree $2 + 3 = 5,$ we conclude that $R(z) - P(z)\cdot Q(z)$ has degree $6$ and is thus nonconstant.

 It now suffices to illustrate an example for which $N = 1$: Take
 \begin{align*} P(z)&=z^2+1, \\ Q(z)&=z^3+2, \\ R(z)&=(z+1)^6 + P(z) \cdot Q(z). \end{align*}
Note that $R(z)$ has degree $6$ and constant term $3,$ so it satisfies the conditions.

 We need to find the solutions to 
 \begin{align*} P(z) \cdot Q(z) &= (z+1)^6 + P(z) \cdot Q(z) \\ 0 &= (z+1)^6. \end{align*}
Clearly, the only distinct complex root is $-1,$ so our answer is $N=\boxed{\textbf{(B)} \: 1}.$



	\item (\textit{AMC 12 12B Fall 2021}) Three identical square sheets of paper each with side length $6$ are stacked on top of each other. The middle sheet is rotated clockwise $30^\circ$ about its center and the top sheet is rotated clockwise $60^\circ$ about its center, resulting in the $24$-sided polygon shown in the figure below. The area of this polygon can be expressed in the form $a-b\sqrt{c}$, where $a$, $b$, and $c$ are positive integers, and $c$ is not divisible by the square of any prime. What is $a+b+c$?


 
                    \begin{figure}[H]
                        \centering
                        \includegraphics[width=0.5\textwidth]{Problem/AMC_12_12B_Fall_2021_15_Problem_Figure_1.png}
                    \end{figure}
                    
$(\textbf{A})\: 75\qquad(\textbf{B}) \: 93\qquad(\textbf{C}) \: 96\qquad(\textbf{D}) \: 129\qquad(\textbf{E}) \: 147$

 \textbf{Answer: }E



\textbf{Solution 1}


                    \begin{figure}[H]
                        \centering
                        \includegraphics[width=0.5\textwidth]{Solution/AMC_12_12B_Fall_2021_15_Solution_Figure_1.png}
                    \end{figure}
                    The $24$-sided polygon is made out of $24$ shapes like $\triangle ABC$. Then $\angle BAC=360^\circ/24=15^\circ$, and $\angle EAC = 45^\circ$, so  $\angle{EAB} = 30^{\circ}$. Then $EB=AE\tan 30^\circ = \sqrt{3}$; therefore $BC=EC-EB=3-\sqrt{3}$. Thus
 \[[ABC] = \frac{BC}{EC}\cdot [ACE] = \frac{3-\sqrt{3}}{3}\cdot \frac 92\]and the required area is $24\cdot[ABC] =108-36\sqrt{3}$. Finally $108+36+3=\boxed{(\textbf{E})\ 147}$. ~lopkiloinm

 Note: Drop an altitude from $A$ to $\overline{BC}$ to construct point $E$. This creates right triangles. ~erringbubble

 



\textbf{Solution 2}

As shown in \href{https://artofproblemsolving.com/wiki/index.php?title=File:2021\_AMC\_12B\_(Nov)\_Problem\_15,\_sol.png}{Image:2021\_AMC\_12B\_(Nov)\_Problem\_15,\_sol.png}, all 12 vertices of three squares form a regular dodecagon (12-gon).Denote by $O$ the center of this dodecagon.

 Hence, $\angle AOB = \frac{360^\circ}{12} = 30^\circ$.

 Because the length of a side of a square is 6, $AO = 3 \sqrt{2}$.

 Hence, $AB = 2 AO \sin \frac{\angle AOB}{2} = 3 \left( \sqrt{3} - 1 \right)$.

 We notice that $\angle MAB = \angle MBA = 30^\circ$.Hence, $AM = \frac{AB}{2\cos \angle MAB} = 3 - \sqrt{3}$.

 Therefore, the area of the region that three squares cover is
 \begin{align*} & {\textrm Area} \ ABCDEFGHIJKL - 12 {\textrm Area} \ \triangle MAB \\ & = 12 {\textrm Area} \ \triangle OAB - 12 {\textrm Area} \ \triangle MAB \\ & = 12 \cdot \frac{1}{2} OA \cdot OB \sin \angle AOB - 12 \cdot \frac{1}{2} MA \cdot MB \sin \angle AMB \\ & = 6 OA^2 \sin \angle AOB - 6 MA^2 \sin \angle AMB \\ & = 108 - 36 \sqrt{3} . \end{align*}


 Therefore, the answer is $\boxed{\textbf{(E) }147}$.

 



\textbf{Solution 3 (complex number & coordinate geometry)}


                    \begin{figure}[H]
                        \centering
                        \includegraphics[width=0.5\textwidth]{Solution/AMC_12_12B_Fall_2021_15_Solution_Figure_2.png}
                    \end{figure}
                    set A = 3+3i , A' , B'  rotate 30 degree from A, B    

 A'= A $\cdot e^i30^\circ  = (3+3i)*(\sqrt{3}/2  + 1/2 i) =( \sqrt{3}/2 - 1/2) + (1/2 + \sqrt{3}/2) i$ 

 line A'B' $\frac{y-(1/2 + \sqrt{3}/2)}{x - ( \sqrt{3}/2 - 1/2)} = Tan(90\circ+30\circ) = -\sqrt{3}$ 

 intersect with line y=3  at point $E_{x} = \sqrt{3}$  , then length $AE = A_{x} - E_{x} = 3- \sqrt{3}$ ,  

 use shoelace or $\triangle OAE$ = 1/2 * AE * AB/2 = 1/2 * $(3- \sqrt{3})$ * 3  

 total area = 24 * $\triangle OAE$  = = 108 - 36 $\sqrt{3}$ the answer is $\boxed{\textbf{(E) }147}$.

 


	\item (\textit{AMC 12 12B Fall 2021}) Suppose $a$, $b$, $c$ are positive integers such that 
 \[a+b+c=23\] and 
 \[\gcd(a,b)+\gcd(b,c)+\gcd(c,a)=9.\] What is the sum of all possible distinct values of $a^2+b^2+c^2$? 


 $\textbf{(A)}\: 259\qquad\textbf{(B)} \: 438\qquad\textbf{(C)} \: 516\qquad\textbf{(D)} \: 625\qquad\textbf{(E)} \: 687$

 \textbf{Answer: }B



\textbf{Solution 1 (Observation)}

Because $a + b + c$ is odd, $a$, $b$, $c$ are either one odd and two evens or three odds.

 $\textbf{Case 1}$: $a$, $b$, $c$ have one odd and two evens.

 Without loss of generality, we assume $a$ is odd and $b$ and $c$ are even.

 Hence, ${\textrm gcd} \left( a , b \right)$ and ${\textrm gcd} \left( a , c \right)$ are odd, and ${\textrm gcd} \left( b , c \right)$ is even.Hence, ${\textrm gcd} \left( a , b \right) + {\textrm gcd} \left( b , c \right) + {\textrm gcd} \left( c , a \right)$ is even. This violates the condition given in the problem.

 Therefore, there is no solution in this case.

 $\textbf{Case 2}$: $a$, $b$, $c$ are all odd.

 In this case, ${\textrm gcd} \left( a , b \right)$, ${\textrm gcd} \left( a , c \right)$, ${\textrm gcd} \left( b , c \right)$ are all odd.

 Without loss of generality, we assume
 \[ {\textrm gcd} \left( a , b \right) \leq {\textrm gcd} \left( b , c \right) \leq {\textrm gcd} \left( c , a \right) . \]$\textbf{Case 2.1}$: ${\textrm gcd} \left( a , b \right) = 1$, ${\textrm gcd} \left( b , c \right) = 1$, ${\textrm gcd} \left( c , a \right) = 7$.

 The only solution is $(a, b, c) = (7, 9, 7)$.

 Hence, $a^2 + b^2 + c^2 = 179$.

 $\textbf{Case 2.2}$: ${\textrm gcd} \left( a , b \right) = 1$, ${\textrm gcd} \left( b , c \right) = 3$, ${\textrm gcd} \left( c , a \right) = 5$.

 The only solution is $(a, b, c) = (5, 3, 15)$.

 Hence, $a^2 + b^2 + c^2 = 259$.

 $\textbf{Case 2.3}$: ${\textrm gcd} \left( a , b \right) = 3$, ${\textrm gcd} \left( b , c \right) = 3$, ${\textrm gcd} \left( c , a \right) = 3$.

 There is no solution in this case.

 Therefore, putting all cases together, the answer is $179 + 259 = \boxed{\textbf{(B)} \: 438}$.

 



\textbf{Solution 2 (Enumeration)}

Let $\gcd(a,b)=x$, $\gcd(b,c)=y$, $\gcd(c,a)=z$. Without the loss of generality, let $x \le y \le z$. We can split this off into cases: 

 $x=1,y=1,z=7$: let $a=7A, c=7C,$ we can try all possibilities of $A$ and $C$ to find that $a=7, b=9, c=7$ is the only solution. 

 $x=1,y=2,z=6$: No solutions. By $y$ and $z$, we know that $a$, $b$, and $c$ have to all be divisible by $2$. Therefore, $x$ cannot be equal to $1$. 

 $x=1,y=3,z=5$: Note that $c$ has to be both a multiple of $3$ and $5$. Therefore, $c$ has to be a multiple of $15$. The only solution for this is $a=5, b=3, c=15$. 

 $x=1,y=4,z=4$: No solutions. By $y$ and $z$, we know that $a$, $b$, and $c$ have to all be divisible by $4$. Therefore, $x$ cannot be equal to $1$. 

 $x=2,y=2,z=5$: No solutions. By $x$ and $y$, we know that $a$, $b$, and $c$ have to all be divisible by $2$. Therefore, $z$ cannot be equal to $5$. 

 $x=2,y=3,z=4$: No solutions. By $x$ and $z$, we know that $a$, $b$, and $c$ have to all be divisible by $2$. Therefore, $y$ cannot be equal to $3$. 

 $x=3,y=3,z=3$: No solutions. As $a$, $b$, and $c$ have to all be divisible by $3$, $a+b+c$ has to be divisible by $3$. This contradicts the sum $a+b+c=23$. 

 Putting these solutions together, we have $(7^2+9^2+7^2)+(5^2+3^2+15^2)=179+259=\boxed{\textbf{(B)} \: 438}$.

 



\textbf{Solution 3}

Since $a+b+c=23$, $\gcd(a,b,c)=23$ or $\gcd(a,b,c)=1$.

 As $\gcd(a,b)+\gcd(b,c)+\gcd(c,a)=9$, it is impossible for $\gcd(a,b,c)=23$, so $\gcd(a,b,c)=1$.

 This means that $\gcd(a,b)$, $\gcd(b,c)$, and $\gcd(c,a)$ must all be coprime. The only possible ways for this to be true are $1+1+7=9$ and $1+3+5=9$.

 Without loss of generality, let $a\le b\le c$. Since $a+b+c=23$, then $a=7, b=7, c=9$ or $a=3, b=5, c=15$.

 $(7^2+7^2+9^2)+(3^2+5^2+15^2)=179+259=\boxed{\textbf{(B)} \: 438}$.

 


	\item (\textit{AMC 12 12B Fall 2021}) A bug starts at a vertex of a grid made of equilateral triangles of side length $1$. At each step the bug moves in one of the $6$ possible directions along the grid lines randomly and independently with equal probability. What is the probability that after $5$ moves the bug never will have been more than $1$ unit away from the starting position?


 $\textbf{(A)}\ \frac{13}{108} \qquad\textbf{(B)}\  \frac{7}{54} \qquad\textbf{(C)}\  \frac{29}{216} \qquad\textbf{(D)}\ \frac{4}{27} \qquad\textbf{(E)}\ \frac{1}{16}$

 \textbf{Answer: }A



\textbf{Solution 1 (Recursion)}

Let $S(n)$ be the number of paths of $n$ moves such that the bug never will have been more than $1$ unit away from the starting position. Clearly, by symmetry, there are two possible states here, the bug being on the center and the bug being on one of the vertices of the unit hexagon around the center. Let $C(n)$ be the number of paths with the aforementioned restriction that end on the center. Let $V(n)$ be the number of paths with the aforementioned restriction that end on a vertex of the surrounding unit hexagon. We have $S(n) = 6C(n-1) + 3V(n-1),$ since from the center, there are $6$ possible points to land to and from a vertex there are $3$ possible points to land to (the two adjacent vertices and the center). We also have $C(n) = V(n-1)$, since to get to the center the bug must have come from a vertex, and $V(n) = 2V(n-1) + 6C(n-1),$ since from a vertex there are two vertices to move to, and from the center there are $6$ vertices to move to. We can construct a recursion table using the base cases $V(1) = 6$ and $C(1) = 0$ and our recursive rules for $C(n)$ and $V(n)$ as follows: 
 \[\begin{array}{c|c|c} n & V(n) & C(n) \\ \hline & & \\ [-2ex] 1 & 6 & 0 \\ 2 & 12 & 6 \\ 3 & 60 & 12 \\ 4 & 192 & 60 \\ \end{array}\]Then, $S(5) = 6C(4) + 3V(4) = 6 \cdot 60 + 3 \cdot 192 = 936,$ and the desired probability is thus $\frac{936}{6^5} = \boxed{\textbf{(A)}\ \frac{13}{108}}.$




\textbf{Solution 2 (Recursion)}

Let $p(i)$ be such probability after $i$ moves. $p(1)=1$, $p(2)=\frac{1}{2}$. Then, $p(3)=\frac{1}{3}\cdot \frac{1}{2}+\frac{1}{6}\cdot 1=\frac{1}{3}$. Then, we can prove the recursive formula 
 \[p(x)=\frac{1}{3}p(x-1)+\frac{1}{6}p(x-2).\] Now, we evaluate $p(5)=\boxed{\textbf{(A)}\ \frac{13}{108}}$.

 



\textbf{Solution 3 (Recursion)}

We use $(n,t)$ to denote the bug's current state. We wish to find $P(0,5)$.

 The first argument $n$ denotes the bug's current position.We use $n = 0$ to denote the bug's starting point.We use $n = 1$ to denote any point whose distance to the bug's starting point is $1$.

 The second argument $t$ denotes the remaining number of moves the bug has.

 For $n = 0$ and $t \geq 1$, we have 
 \[P(0,t) = P(1,t-1).\]For $n = 1$ and $t \geq 1$, we have 
 \[P(1,t) = \frac{1}{6} P(0,t-1) + \frac{1}{3} P(1,t-1).\]For $n \in \left\{ 0 , 1 \right\}$ and $t = 0$, we have 
 \[P(n,0) = 1.\]We solve this recursive equation by using backward induction:
 \[\begin{array}{ll} P(0,1) = 1, & P(1,1) = \frac{1}{2}, \\ [1ex] P(0,2) = \frac{1}{2}, & P(1,2) = \frac{1}{3}, \\ [1ex] P(0,3) = \frac{1}{3}, & P(1,3) = \frac{7}{36}, \\ [1ex] P(0,4) = \frac{7}{36}, & P(1,4) = \frac{13}{108}. \end{array}\]Therefore, the answer is $P(0,5) = P(1,4) = \boxed{\textbf{(A)}\ \frac{13}{108}}$.

 



\textbf{Solution 4 (Generating Function)}

Use a generating function, define $c_{n}\cdot x^{n}$ be $c_{n}$ ways for the destination be $n$ units away from the origin. 

 We conclude that:

 
            \begin{itemize}
        \item If the current point is origin, then we need to multiply by $6x$.
            \end{itemize}
            
            \begin{itemize}
        \item If the current point on vertex of the unit hexagon, then we need to multiply by $x^{-1}+2$, where there is $1$ way to return to the origin and there are two ways to keep distance $1$.
            \end{itemize}
            Now let's start with $p(x)=1$.

 $1$st step:  $p(x)=6x$
$2$nd step:  $p(x)=6x\cdot(x^{-1}+2) = 6 + 12x$
$3$rd step:  $p(x)=6\cdot6x + 12x\cdot(x^{-1}+2) = 12 + 60x$
$4$th step:  $p(x)=12\cdot6x + 60x\cdot(x^{-1}+2) = 60 + 192x$
$5$th step:  $p(x)=60\cdot6x + 192x\cdot(x^{-1}+2) = 192 + 744x$

 So, there are $192+744=936$ ways for the bug never moves more than $1$ unit away from origin. The answer is $\frac{936}{6^5} = \boxed{\textbf{(A)}\ \frac{13}{108}}$.

 



\textbf{Solution 5 (Casework)}

In the following diagram, let $A$ denote the vertex where the bug starts (shown in red) and $B$ denote one of the $6$ adjacent vertices (shown in green).
                    \begin{figure}[H]
                        \centering
                        \includegraphics[width=0.5\textwidth]{Solution/AMC_12_12B_Fall_2021_17_Solution_Figure_1.png}
                    \end{figure}
                    Note that:

 
            \begin{itemize}
        \item If the bug is at $A,$ then the probability that it moves to $B$ next is $1.$
            \end{itemize}
            
            \begin{itemize}
        \item If the bug is at $B,$ then the probability that it moves to $A$ next is $\frac16.$
            \end{itemize}
            
            \begin{itemize}
        \item If the bug is at $B,$ then the probability that it moves to $B$ next is $\frac13.$
            \end{itemize}
            We apply casework to the possible paths of the bug:

 
            \begin{enumerate} 
        \item $A \rightarrow B \rightarrow A \rightarrow B \rightarrow A \rightarrow B$ The probability for this case is $1\cdot\frac16\cdot1\cdot\frac16\cdot1=\frac{1}{36}.$


        \item $A \rightarrow B \rightarrow A \rightarrow B \rightarrow B \rightarrow A$ The probability for this case is $1\cdot\frac16\cdot1\cdot\frac13\cdot\frac16=\frac{1}{108}.$


        \item $A \rightarrow B \rightarrow B \rightarrow A \rightarrow B \rightarrow A$ The probability for this case is $1\cdot\frac13\cdot\frac16\cdot1\cdot\frac16=\frac{1}{108}.$


        \item $A \rightarrow B \rightarrow A \rightarrow B \rightarrow B \rightarrow B$ The probability for this case is $1\cdot\frac16\cdot1\cdot\frac13\cdot\frac13=\frac{1}{54}.$


        \item $A \rightarrow B \rightarrow B \rightarrow A \rightarrow B \rightarrow B$ The probability for this case is $1\cdot\frac13\cdot\frac16\cdot1\cdot\frac13=\frac{1}{54}.$


        \item $A \rightarrow B \rightarrow B \rightarrow B \rightarrow A \rightarrow B$ The probability for this case is $1\cdot\frac13\cdot\frac13\cdot\frac16\cdot1=\frac{1}{54}.$


        \item $A \rightarrow B \rightarrow B \rightarrow B \rightarrow B \rightarrow A$ The probability for this case is $1\cdot\frac13\cdot\frac13\cdot\frac13\cdot\frac16=\frac{1}{162}.$


        \item $A \rightarrow B \rightarrow B \rightarrow B \rightarrow B \rightarrow B$ The probability for this case is $1\cdot\frac13\cdot\frac13\cdot\frac13\cdot\frac13=\frac{1}{81}.$


            \end{enumerate}
            Together, the answer is 
 \[\frac{1}{36}+\frac{1}{108}+\frac{1}{108}+\frac{1}{54}+\frac{1}{54}+\frac{1}{54}+\frac{1}{162}+\frac{1}{81}=\boxed{\textbf{(A)}\ \frac{13}{108}}.\]~MRENTHUSIASM

 



\textbf{Solution 6 (Markov Chain)}

We can use a \href{https://artofproblemsolving.com/wiki/index.php?title=Markov\_Chains}{Markov chain} to represent the different states of the bug. 

 
                    \begin{figure}[H]
                        \centering
                        \includegraphics[width=0.5\textwidth]{Solution/AMC_12_12B_Fall_2021_17_Solution_Figure_2.png}
                    \end{figure}
                    Let "center" (C) represent the starting point, "1 away" (O) represent a point on the grid 1 move and 1 unit away from the starting point, and "2+ away" (T) represent the case where the bug has ventured more than 1 unit away from its starting point.

 Notice that the bug must move to the 1 away state from the center. From there, it moves back to the center with a $\frac{1}{6}$ probability, to a point 2 units away from the center with $\frac{1}{2}$ probability, and to another point 1 unit away with $\frac{1}{3}$ probability.

 We are trying to find the probability that the bug does not reach the 2+ away state in 5 moves. To accomplish this, we can find the complement by listing out the cases that the bug does reach the 2+ away state in a specified number of moves.

 Clearly, the bug cannot reach the 2+ away (T) state in 0 moves.

 In 1 move, the bug can only reach the 1 away (O) state, so it cannot reach the T state in 1 move either.

 In 2 moves, the bug can move to the O state followed by the T state with probability $1*\frac{1}{2} = \frac{1}{2}$.

 In 3 moves, the bug can follow the pattern OOT with probability $1*\frac{1}{3}*\frac{1}{2} = \frac{1}{6}$.

 In 4 moves, the bug can follow the paths of OOOT or OCOT, with probabilities $\frac{1}{18}$ and $\frac{1}{12}$ respectively.

 In 5 moves, the bug can follow the paths of OOOOT, OCOOT, or OOCOT, with probabilities $\frac{1}{54}$, $\frac{1}{36}$, and $\frac{1}{36}$. 

 Adding up these probabilities, we obtain $\frac{95}{108}$. Finding the complement of this, we subtract from 1 to get the answer of $\boxed{\textbf{(A)}\ \frac{13}{108}}$.$\newline$~diyarv

 


	\item (\textit{AMC 12 12B Fall 2021}) Set $u_0 = \frac{1}{4}$, and for $k \ge 0$ let $u_{k+1}$ be determined by the recurrence 
 \[u_{k+1} = 2u_k - 2u_k^2.\]


 This sequence tends to a limit; call it $L$. What is the least value of $k$ such that 
 \[|u_k-L| \le \frac{1}{2^{1000}}?\]
$\textbf{(A)}\: 10\qquad\textbf{(B)}\: 87\qquad\textbf{(C)}\: 123\qquad\textbf{(D)}\: 329\qquad\textbf{(E)}\: 401$

 \textbf{Answer: }A



\textbf{Solution 1}

Note that terms of the sequence $(u_k)$ lie in the interval $\left(0,\frac12\right)$ and are strictly increasing.

 Since the sequence $(u_k)$ tends to the limit $L,$ we set $u_{k+1}=u_k=L>0.$

 The given equation becomes 
 \[L=2L-2L^2,\] from which $L=\frac12.$

 The given inequality becomes 
 \[\frac12-\frac{1}{2^{1000}} \leq u_k \leq \frac12+\frac{1}{2^{1000}},\] and we only need to consider $\frac12-\frac{1}{2^{1000}} \leq u_k.$

 We have
 \begin{alignat*}{8} u\_0 &= \phantom{1}\frac14 &&= \frac{2\textasciicircum{}1-1}{2\textasciicircum{}2}, \\ u\_1 &= \phantom{1}\frac38 &&= \frac{2\textasciicircum{}2-1}{2\textasciicircum{}3}, \\ u\_2 &= \ \frac{15}{32} &&= \frac{2\textasciicircum{}4-1}{2\textasciicircum{}5}, \\ u\_3 &= \frac{255}{512} &&= \frac{2\textasciicircum{}8-1}{2\textasciicircum{}9}, \\ & \phantom{1111} \vdots \end{alignat*}By induction, it can be proven that 
 \[u_k=\frac{2^{2^k}-1}{2^{2^k+1}}=\frac12-\frac{1}{2^{2^k+1}}.\]We substitute this into the inequality, then solve for $k:$
 \begin{align*} \frac12-\frac{1}{2^{1000}} &\leq \frac12-\frac{1}{2^{2^k+1}} \\ -\frac{1}{2^{1000}} &\leq -\frac{1}{2^{2^k+1}} \\ 2^{1000} &\leq 2^{2^k+1} \\ 1000 &\leq 2^k+1. \end{align*}
Since $2^9+1 \leq 1000 \leq 2^{10}+1,$ the least such value of $k$ is $\boxed{\textbf{(A)}\: 10}.$




\textbf{Solution 2}

If we list out the first few values of $k$, we get the series $\frac{1}{4}, \frac{3}{8}, \frac{15}{32}, \frac{255}{512}$, which always seems to be a negative power of $2$ away from $\frac{1}{2}$. We can test this out by setting $u_k=\frac{1}{2}-\frac{1}{2^{n_k}}$, where $n_0=2$.

 Now, we get
 \begin{align*} u_{k+1} &= 2\cdot\left(\frac{1}{2}-\frac{1}{2^{n_{k}}}\right)-2\cdot\left(\frac{1}{2}-\frac{1}{2^{n_{k}}}\right)^2 \\ &= 1-\frac{1}{2^{n_k - 1}}-2\cdot\left(\frac{1}{4}-\frac{1}{2^{n_k}}+\frac{1}{2^{2 \cdot n_k}}\right)\\ &= 1-\frac{1}{2^{n_k - 1}}-\frac{1}{2}+\frac{1}{2^{n_k-1}}-\frac{1}{2^{2 \cdot n_k-1}} \\ &= \frac{1}{2}-\frac{1}{2^{2 \cdot n_k-1}}. \end{align*}
This means that this series approaches $\frac{1}{2}$, as the second term is decreasing. In addition, we find that $n_{k+1}=2 \cdot n_k-1$. 

 We claim that $n_k = 2^k+1$, which can be proven by induction:

 <u>\textbf{Base Case}</u>

 We have $n_0=2=2^0+1$.

 <u>\textbf{Induction Step}</u>

 Assuming that the claim is true, we have $n_{k+1}=2 \cdot (2^k+1)-1=2^{k+1}+1$.

 It follows that $n_{10}=2^{10}+1>1000$ and $n_9=2^9+1<1000$. Therefore, the least value of $k$ would be $\boxed{\textbf{(A)}\: 10}$.

 



\textbf{Solution 3}

We are given $u_{k+1} = 2u_k - 2{u_k}^2$. Multiply this equation by $2$ and subtract $1$ from both sides. The equations can then be written nicely as $2u_{k+1} - 1 = -(2u_k-1)^2$. Let $v_k = 2u_k - 1$ so that $v_{k+1} = -(v_k)^2$.

 Clearly, $v_0 = 2u_0 - 1 = -\frac{1}{2}$. Since the magnitude of $v_0$ is less than $1$ and because our recursive relation for $v_k$ squares the previous term (and negates it), we see that as $k \rightarrow \infty, 2u_k - 1 = v_k \rightarrow 0$. This means $u_k \rightarrow \frac{1}{2}$, so $L = \frac{1}{2}$.

 Isolating $u_k$ in our relation $2u_k - 1 = v_k$ gives us $u_k = \frac{v_k + 1}{2}$.Substituting into the inequality, we have $\left|\frac{v_k + 1}{2}-\frac{1}{2}\right| \le \frac{1}{2^{1000}}$. Rewriting this, we get $|v_k| \le \frac{1}{2^{999}}$.

 The sequence $\{v_k\}$ is much easier to handle because of its simple recursive relation. Writing out a few terms shows that $|v_k| = \frac{1}{2^{2^k}}$. Now it just comes down to having $2^k > 999$, so $k = \boxed{\textbf{(A)}\: 10}$.

 


	\item (\textit{AMC 12 12B Fall 2021}) Regular polygons with $5,6,7,$ and $8$ sides are inscribed in the same circle. No two of the polygons share a vertex, and no three of their sides intersect at a common point. At how many points inside the circle do two of their sides intersect?


 $(\textbf{A})\: 52\qquad(\textbf{B}) \: 56\qquad(\textbf{C}) \: 60\qquad(\textbf{D}) \: 64\qquad(\textbf{E}) \: 68$

 \textbf{Answer: }E



\textbf{Solution}

Imagine we have $2$ regular polygons with $m$ and $n$ sides and $m>n$ inscribed in a circle without sharing a vertex. We see that each side of the polygon with $n$ sides (the polygon with fewer sides) will be intersected twice.(We can see this because to have a vertex of the $m$-gon on an arc subtended by a side of the $n$-gon, there will be one intersection to “enter” the arc and one to “exit” the arc. ~KingRavi)

 This means that we will end up with $2$ times the number of sides in the polygon with fewer sides.

 If we have polygons with $5,$ $6,$ $7,$ and $8$ sides, we need to consider each possible pair of polygons and count their intersections. 

 Throughout $6$ of these pairs, the $5$-sided polygon has the least number of sides $3$ times, the $6$-sided polygon has the least number of sides $2$ times, and the $7$-sided polygon has the least number of sides $1$ time.

 Therefore the number of intersections is $2\cdot(3\cdot5+2\cdot6+1\cdot7)=\boxed{(\textbf{E}) \:68}$.

 \textbf{Remark} For regular polygons with $5,6,7,$ and $8$ sides, the $68$ points of intersection inside the circle are shown below:
                    \begin{figure}[H]
                        \centering
                        \includegraphics[width=0.5\textwidth]{Solution/AMC_12_12B_Fall_2021_19_Solution_Figure_1.png}
                    \end{figure}
                    ~MRENTHUSIASM

 


	\item (\textit{AMC 12 12B Fall 2021}) A cube is constructed from $4$ white unit cubes and $4$ blue unit cubes. How many different ways are there to construct the $2 \times 2 \times 2$ cube using these smaller cubes? (Two constructions are considered the same if one can be rotated to match the other.)


 $(\textbf{A})\: 7\qquad(\textbf{B}) \: 8\qquad(\textbf{C}) \: 9\qquad(\textbf{D}) \: 10\qquad(\textbf{E}) \: 11$

 \textbf{Answer: }A



\textbf{Solution 1 (Graph Theory)}

This problem is about the relationships between the white unit cubes and the blue unit cubes, which can be solved by \href{https://en.wikipedia.org/wiki/Graph\_theory}{Graph Theory}. We use a \href{https://en.wikipedia.org/wiki/Planar\_graph}{Planar Graph} to represent the larger cube. Each vertex of the planar graph represents a unit cube. Each edge of the planar graph represents a shared face between $2$ neighboring unit cubes. Each face of the planar graph represents a face of the larger cube.

 Now the problem becomes a \href{https://en.wikipedia.org/wiki/Graph\_coloring}{Graph Coloring} problem of how many ways to assign $4$ vertices blue and $4$ vertices white with \href{https://en.wikipedia.org/wiki/Homeomorphism}{Topological Equivalence}. For example, in Figure $(1)$, as long as the $4$ blue vertices belong to the same planar graph face, the different planar graphs are considered to be topological equivalent by rotating the larger cube.

 
                    \begin{figure}[H]
                        \centering
                        \includegraphics[width=0.5\textwidth]{Solution/AMC_12_12B_Fall_2021_20_Solution_Figure_1.jpg}
                    \end{figure}
                    

 Here is how the $4$ blue unit cubes are arranged: 

 In Figure $(1)$: $4$ blue unit cubes are on the same layer (horizontal or vertical).

 In Figure $(2)$: $4$ blue unit cubes are in $T$ shape.

 In Figure $(3)$ and $(4)$: $4$ blue unit cubes are in $S$ shape.

 In Figure $(5)$: $3$ blue unit cubes are in $L$ shape, and the other is isolated without a shared face.

 In Figure $(6)$: $2$ pairs of neighboring blue unit cubes are isolated from each other without a shared face.

 In Figure $(7)$: $4$ blue unit cubes are isolated from each other without a shared face.

 So the answer is $\boxed{\textbf{(A)}\ 7}$




\textbf{Solution 2 (Casework, Counting Down)}

Let’s split the cube into two layers; a bottom and top. Note that there must be four of each color, so however many number of one color are in the bottom, there will be four minus that number of the color on the top. We do casework on the color distribution of the bottom layer.

 \textbf{Case 1:} 4, 0

 In this case, there is only one possibility for the top layer - all of the other color - $\binom{4}{4}$. Therefore there is 1 construction from this case.

 \textbf{Case 2:} 3, 1

 In this case, the top layer has four possibilities, because there are four different ways to arrange it so that it also has a 3, 1 color distribution - $\binom{4}{3}$. Therefore there are 4 constructions from this case.

 \textbf{Case 3:} 2, 2

 In this case, the top layer has six possibilities of arrangement - $\binom{4}{2}$. However, having adjacent colors one way can be rotated to having adjacent colors any other way, so there is only one construction for the adjacent colors subcase and similarly, only one for the diagonal color subcase. Therefore the total number of constructions for this case is 2.

 The total number of constructions for the cube is thus $1+4+2=7=\boxed{A}$




\textbf{Solution 3 (Casework, Counting Up)}

Divide the $2 \times 2 \times 2$ cube into two layers, say, front and back. Any possible construction can be rotated such that the front layer has the same or greater number of white cubes than blue cubes, so we only need to count the number of cases given that is true.  

 
            \begin{enumerate}
        \item \textbf{Case 1: Each layer contains 2 cubes of each color.} Note that we only need to consider the configuration of the white cubes because all the other cubes will be blue cubes. There are 2 ways that the 2 white cubes in each layer can be arranged: adjacent or diagonal to each other.<ol><li>\textbf{Case 1.1: Both layers have 2 white cubes adjacent to each other.} Rotate the cube such that there are white cubes along the top edge of the front layer. Now, the white cubes in the back layer can be along the top, bottom, right, or left edges. So, case 1.1 results in $4$ constructions.</li><li>\textbf{Case 1.2: One layer has 2 white cubes adjacent to each other, and the other has 2 white cubes diagonal from each other.} Rotate the cube such that there are white cubes along the top of the front layer. The white cubes in the back layer can be at the top-left and bottom-right or at the top-right and bottom-left. If we rotate the former case by 90 degrees clockwise into the page, it becomes the same as the latter case. So, case 1.2 results in $1$ additional construction.</li><li>\textbf{Case 1.3: Both layers have white cubes diagonal from each other.} Rotate the cube such that there is a white cube at the top-left and bottom-right of the front layer. The back layer could also have white cubes at the top-left and bottom-right, but this is the same as case 1.1 with the white cubes in the back layer along the bottom edge. Alternatively, the back layer could have white cubes at the top-right and the bottom-left. This is a distinct case. So, case 1.3 results in $1$ additional construction.</li><li>So, case 1 results in $4+1+1=6$ distinct constructions.</li></ol>
        \item \textbf{Case 2: The front layer contains 3 white cubes, and the back layer contains one white cube.} In this case, unless the sole white and sole blue cubes in the front and back layers are on opposite corners of the $2 \times 2 \times 2$ cube, then that cube can be split into left and right layers with 2 cubes of each color in each layer. Those constructions were counted in case 1, so case 2 results in $1$ additional construction.
        \item \textbf{Case 3: The front layer contains 4 white cubes, and the back layer contains no white cubes.} If we split this construction into its left and right layers, then each layer will have 2 cubes of each color. So, this construction was counted in case 1, and case 3 results in $0$ additional constructions.
            \end{enumerate}
            Therefore, our answer is $6+1+0=\boxed{\textbf{(A)}\ 7}$.

 



\textbf{Solution 4 (Burnside Lemma)}

Burnside lemma is used to counting number of orbit where the element on the same orbit can be achieved by the defined operator, naming rotation, reflection and etc. 

 The facts for Burnside lemma are 

 1. the sum of stabilizer on the same orbit equals to the \# of operators; 

 2. the sum of stabilizer can be counted as $fix(g)$ 

 3. the sum of the $fix(g)/|G|$ equals the \# of orbit. 

 
Let's start with defining the operator for a cube, 

 1. $\textbf{e (identity)}$

 For identity, there are $\frac{8!}{4!4!} = 70$ 

 
2. ${\textbf r^{1}, r^{2}, r^{3}}$ to be the rotation axis along three pair of opposite face, 

 each contains $r^{i}_{90}, r^{i}_{180}, r^{i}_{270}$ where $i= 1, 2, 3$
$fix(r^{i}_{90}) = fix(r^{i}_{270}) = 2\cdot1 = 2$
$fix(r^{i}_{180}) = \frac{4!}{2!\cdot2!} = 6$

 therefore $fix(\textbf r^{i}) = 2+2+6 = 10$, and $fix(\textbf r^{1})+fix(\textbf r^{2})+fix(\textbf r^{3}) = 30$

 
3. ${\textbf r^{4}, r^{5}, r^{6}, r^{7}}$ to the rotation axis along four cube diagonals. 

 each contains $r^{i}_{120}, r^{i}_{240}$   where $i= 4, 5, 6, 7$
$fix(r^{i}_{120}) = fix(r^{i}_{240}) = 2\cdot1\cdot2\cdot1 = 4$

 therefore $fix(\textbf r^{i}) = 4+4 = 8$, and $fix(\textbf r^{4})+fix(\textbf r^{5})+fix(\textbf r^{6})+fix(\textbf r^{7}) = 32$

 


 4. ${\textbf r^{8}, r^{9}, r^{10}, r^{11}, r^{12}, r^{13}}$ to be the rotation axis along 6 pairs of diagonally opposite sides

 each contains $r^{i}_{180}$   where $i= 8, 9, 10, 11, 12, 13$
$fix(r^{i}_{180}) = \frac{4!}{2!\cdot2!} = 6$

 therefore $fix(\textbf r^{8})+fix(\textbf r^{9})+fix(\textbf r^{10})+fix(\textbf r^{11})+fix(\textbf r^{12})+fix(\textbf r^{13}) = 36$

 
5. The total number of operators are 

 $|G| = 1 + 3\cdot3 + 4\cdot2 + 6\cdot1 = 24$

 Based on 1, 2, 3, 4 the total number of stabilizers is $70 + 30 + 32 + 36 = 168$

 therefore the number of orbits $= \frac{168}{G=24} = \boxed{7}$




\textbf{Solution 5 (Quick 'n Easy but not rigorous)}

Since rotations of a single pattern are considered indistinguishable, we can assume that the forward upper right corner of the 2-by-2-by-2 cube is a blue cube (since we can always rotate the big cube to place a blue cube in that spot).

 Once we've assigned this cube to be blue, we note that 3 1-by-1-by-1 cubes share a side with it, 3 1-by-1-by-1 cubes share a corner with it, and 1 1-by-1-by-1 cube does not touch the assigned cube at all, from the perspective of someone who can only see the cube's faces. We'll call the first 3 "adjacent", the second 3 "cornering", and the last one "opposite."

 We can use a little bit of intuition to confirm that due to the rotation condition, we should treat all adjacents as indistinguishable, all cornerings as indistinguishable, and of course the opposite one is unique from all the others. Thus, we can list out like so (keeping in mind that there are 3 adjacents, 3 cornerings, and 1 opposite, and that we're choosing the positions of the remaining 3 blue cubes):

 OAA, OAC, OCC, CCC, CAA, CCA, AAA.

 This gives the answer to be $\boxed{\textbf{(A)}\ 7}$.

 


	\item (\textit{AMC 12 12B Fall 2021}) For real numbers $x$, let 

 \[P(x)=1+\cos(x)+i\sin(x)-\cos(2x)-i\sin(2x)+\cos(3x)+i\sin(3x)\]
where $i = \sqrt{-1}$. For how many values of $x$ with $0\leq x<2\pi$ does 

 \[P(x)=0?\]
$\textbf{(A)}\ 0 \qquad\textbf{(B)}\  1 \qquad\textbf{(C)}\  2 \qquad\textbf{(D)}\ 3 \qquad\textbf{(E)}\ 4$

 \textbf{Answer: }A



\textbf{Solution 1}

Let $a=\cos(x)+i\sin(x)$. Now $P(a)=1+a-a^2+a^3$. $P(-1)=-2$ and $P(0)=1$ so there is a real root $a_1$ between $-1$ and $0$. The other $a$'s must be complex conjugates since all coefficients of the polynomial are real. The magnitude of those complex $a$'s squared is $-\frac{1}{a_1}$, which is greater than $1$. If $x$ is real number then $a$ must have magnitude of $1$, but none of the solutions for $a$ have magnitude of $1$, so the answer is $\boxed{\textbf{(A)}\ 0 }$ ~lopkiloinm

 



\textbf{Solution 2}

For $\textrm{Im}(P(x))=0$, we get 
 \[\sin(2x)=\sin(x)+\sin(3x)=2\sin(2x)\cos(x)\] So either $\sin(2x)=0$, i.e. $x\in\{0,\pi\}$ or $\cos(x)=\frac 12$, i.e. $x\in \{\pi/3, 5\pi/3\}$.

 For none of these values do we get $\textrm{Re}(P(x))=0$.

 Therefore, the answer is $\boxed{\textbf{(A) }0}$.

 



\textbf{Solution 3}

We have
 \begin{align*} P \left( x \right) & = 1 + e^{ix} - e^{i 2x} + e^{i 3x} . \end{align*}


 Denote $y = e^{i x}$. Hence, this problem asks us to find the number of $y$ with $| y| = 1$ that satisfy
 \[ 1 + y - y^2 + y^3 = 0 . \hspace{1cm} (1) \]

 Taking imaginary part of both sides, we have
 \begin{align*} 0 & = {\textrm Im} \ \left( 1 + y - y^2 + y^3 \right) \\ & = \frac{1}{2i} \left( y - \bar y - y^2 + \bar y^2 + y^3 - \bar y^3 \right) \\ & = \frac{y - \bar y}{2i} \left( 1 - y - \bar y + y^2 + y \bar y + \bar y^2 \right) \\ & = {\textrm Im} \ y \left( 1 - \left( y + \bar y \right) + \left( y + \bar y \right)^2 - y \bar y \right) \\ & = {\textrm Im} \ y \left( 1 - 2 {\textrm Re} \ y + 4 \left( {\textrm Re} \ y \right)^2 - |y|^2 \right) \\ & = {\textrm Im} \ y \left( 1 - 2 {\textrm Re} \ y + 4 \left( {\textrm Re} \ y \right)^2 - 1 \right) \\ & = 2 {\textrm Im} \ y \cdot {\textrm Re} \ y \left( 2 {\textrm Re} \ y - 1 \right) \\ \end{align*}
The sixth equality follows from the property that $|y| = 1$.

 Therefore, we have either ${\textrm Re} \ y = 0$ or ${\textrm Im} \ y = 0$ or $2 {\textrm Re} \ y - 1 = 0$.

 Case 1: ${\textrm Re} \ y = 0$.

 Because $|y| = 1$, $y = \pm i$.

 However, these solutions fail to satisfy Equation (1).

 Therefore, there is no solution in this case.

 Case 2: ${\textrm Im} \ y = 0$.

 Because $|y| = 1$, $y = \pm 1$.

 However, these solutions fail to satisfy Equation (1).

 Therefore, there is no solution in this case.

 Case 3: $2 {\textrm Re} \ y - 1 = 0$.

 Because $|y| = 1$, $y = \frac{1}{2} \pm \frac{\sqrt{3}}{2} i = e^{i \pm \frac{\pi}{3}}$.

 However, these solutions fail to satisfy Equation (1).

 Therefore, there is no solution in this case.

 All cases above imply that there is no solution in this problem.

 Therefore, the answer is $\boxed{\textbf{(A) }0}$.

 



\textbf{Solution 4}

Let $a=\cos(x)+i\sin(x)$, so by De Moivre $P(x)=a^3-a^2+a+1$. The problem essentially asks for the number of real roots of $P$ which lie on the complex unit circle.Let $|r|=1$ be a root of $P$, and note that we can't have $r^3-r^2+r=0$, else $P(r)=0$. Thus, suppose henceforth that $r^3-r^2+r \neq 0$. We then have $r^3-r^2+r=r^2(r+\frac{1}{r}-1)=a^2(2\mathrm{Re}(r)-1)$, hence the argument of $r^3-r^2+r$ is either the argument of $a^2$, or the argument of $-a^2$. Since $r^3-r^2+r=-1$ is real, it follows that $a^2=\pm 1 \implies a \in \{1,i,-1,-i\}$. Now, we can check all of these values and find that none of them work, yielding an answer of $\boxed{\textbf{(A) }0}$.

 -IAmTheHazard

 



\textbf{Solution 5 (Geometry)}

$P(x)$ can be written equivalently as $P(x) = 1 + cis(x) - cis(2x) + cis(3x).$ Thus, we aim to find $x$ such that the sum of the vectors $cis(x)$, $-cis(2x)$, and $cis(3x)$ is -1. Notice that $cis(x)$, $-cis(2x)$, $cis(3x)$ all lie on the unit circle in the complex plane, and the vector $cis(x) + cis(3x)$ is collinear with $-cis(2x).$ Since $|-cis(2x)| = 1$ and we want the three vectors to sum to -1, we either have $-cis(2x) = 1$ and $cis(x) + cis(3x) = -2$, or $-cis(2x) = -1$ and $cis(x) + cis(3x) = 0.$ If the first condition is true, $cis(x) = cis(3x)=-1.$ This will imply that $x= \pi.$ But then $-cis(2x) = -1$, which violates the condition. Similarly, we can show that the second condition cannot be met either. Thus $P(x)$ does not have any solutions on the interval $[0, 2\pi].$ Therefore, the answer is $\boxed{\textbf{(A) }0}$.

 -mathy\_mathema

 



\textbf{Solution 6 (Trig)}


 \[P(x) = 1 + \cos{x} - \cos2x + \cos3x + i(\sin{x} - \sin2x + \sin3x)\]In other words, $1 + \cos{x} - \cos2x + \cos3x = 0$ and $\sin{x} - \sin2x + \sin3x = 0$. Using trigonometric identities, we know that
 \begin{align*}     1 + \cos{x} - \cos2x + \cos3x     &= 0 \\     &= 1 + \cos{x} - (2\cos^2x - 1) + (4\cos^3x - 3\cos{x}) \\     &= 4\cos^3x - 2\cos^2x - 2\cos{x} + 2 \\     &= 2\cos^3x - \cos^2x - \cos{x} + 1 \\     \sin{x} - \sin2x + \sin3x     &= 0 \\     &= \sin{x} - 2\sin{x}\cos{x} + 3\sin{x} - 4\sin^3x \\     &= 4\sin^3x - 4\sin{x} + 2\sin{x}\cos{x} \\     &= \sin{x}(2\sin^2x - 2 + \cos{x}) \\     &= \sin{x}(-2\cos^2x + \cos{x}). \end{align*}
$1 + \cos{x} - \cos2x + \cos3x = 0$ must satisfy when $\sin{x} = 0$, $\cos{x} = 0$ or $\cos{x} = \frac{1}{2}$. The cases when $\cos{x} = 0$ and $\cos{x} = \frac{1}{2}$ does not work. Moreover, for the case $\sin{x} = 0$ to work, $x = 0, \pi, 2\pi$. However, none of the cases work. Therefore, $\boxed{\textbf{(A)}\ 0}$ is the answer.

 


	\item (\textit{AMC 12 12B Fall 2021}) Right triangle $ABC$ has side lengths $BC=6$, $AC=8$, and $AB=10$. A circle centered at $O$ is tangent to line $BC$ at $B$ and passes through $A$. A circle centered at $P$ is tangent to line $AC$ at $A$ and passes through $B$. What is $OP$?


 $\textbf{(A)}\ \frac{23}{8} \qquad\textbf{(B)}\  \frac{29}{10} \qquad\textbf{(C)}\  \frac{35}{12} \qquad\textbf{(D)}\ \frac{73}{25} \qquad\textbf{(E)}\ 3$

 \textbf{Answer: }C



\textbf{Solution 1}

Let $M$ be the midpoint of $AB$; so $BM=AM=5$. Let $D$ be the point such that $ABCD$ is a rectangle. Then $MO\perp AB$ and $MP\perp AB$. Let $\theta = \angle BAC$; so $\tan\theta = \frac 68 = \frac 34$. Then
 \[OP=MP-MO=AM\cot\theta - BM\tan\theta = 5(\frac 43 - \frac 34) = \boxed{\textbf{(C)}\ \frac{35}{12}}.\]




\textbf{Solution 2}

This one uses the same diagram as Solution 1, except we draw $BP$. After doing angle chasing we find $\triangle BPM \sim \triangle BAC$ and $\frac{BM}{BC} = \frac{PM}{AC}$, resulting in $PM = \frac{20}{3}$.

 We also find that $\triangle BOM \sim \triangle ABC$ and $\frac{BM}{AC} = \frac{OM}{BC}$, resulting in $OM = \frac{15}{4}$. $OP = PM - OM = \frac{35}{12}$.

 -ThisUsernameIsTaken

 



\textbf{Solution 3 (Analytic Geometry)}

In a Cartesian plane, let $C, B,$ and $A$ be $(0,0),(0,6),(8,0)$ respectively.

 By analyzing the behaviors of the two circles, we set $O$ to be $(a,6)$ and $P$ be $(b,8)$.

 Hence derive the two equations:

 $(x-a)^2+(y-6)^2=a^2$
$(x-8)^2+(y-b)^2=b^2$

 
Considering the coordinates of $A$ and $B$ for the two equations respectively, we get:

 $(8-a)^2+(0-6)^2=a^2$
$(0-8)^2+(6-b)^2=b^2$

 Solve to get $a=\frac{25}{4}$ and $b=\frac{25}{3}$

 
Through using the distance formula,

 $OP=\sqrt{(8-\frac{25}{4})^2+(\frac{25}{3}-6)^2}= \boxed{\textbf{(C)}\ \frac{35}{12}}$.

 



\textbf{Solution 4}

Because the circle with center $O$ passes through points $A$ and $B$ and is tangent to line $BC$ at point $B$, $O$ is on the perpendicular bisector of segment $AB$ and $OB \perp BC$.

 Because the circle with center $P$ passes through points $A$ and $B$ and is tangent to line $AC$ at point $A$, $P$ is on the perpendicular bisector of segment $AB$ and $PA \perp AC$.

 Let lines $OB$ and $AP$ intersect at point $D$.Hence, $ACBD$ is a rectangle.

 Denote by $M$ the midpoint of segment $AB$. Hence, $BM = \frac{AB}{2} = 5$.Because $O$ and $P$ are on the perpendicular bisector of segment $AB$, points $M$, $O$, $P$ are collinear with $OM \perp AB$.

 We have $\triangle MOB \sim \triangle CBA$.Hence, $\frac{BO}{AB} = \frac{BM}{AC}$.Hence, $BO = \frac{25}{4}$.Hence, $OD = BD - BO = \frac{7}{4}$.

 We have $\triangle DOP \sim \triangle CBA$.Hence, $\frac{OP}{BA} = \frac{DO}{CB}$.Therefore, $OP = \frac{35}{12}$.

 Therefore, the answer is $\boxed{\textbf{(C) }\frac{35}{12}}$.

 



\textbf{Solution 5}

Let $C$ be the origin, making $B=(0,6)$ and $A=(8,0)$. Let $D$ be the midpoint of $AB$; $D=(4,3)$.

 Notice that both $O$ and $P$ must be on the perpendicular bisector $l$ of $AB$. The slope of $AB$ is $-\dfrac{3}{4}$, making the $l$'s slope be $\dfrac{4}{3}$. Since $l$ passes through $D$, the equation for $l$ becomes

 
 \[y-3=\dfrac{4}{3} (x-4),\]

 using the slope intersect form. Since $OB$ is perpendicular to $AC$ and $AP$ is perpendicular to $AC$ (cause of tangencies), the $y$-coordinate for $O$ is $6$ and the $x$-coordinate for $P$ is $8$. Plugging these numbers in the equation for $l$ gives $O=\left( \dfrac{25}{4}, 6 \right)$ and $P=\left( 8, \dfrac{25}{3} \right)$. Thus,

 
 \[OP=\sqrt{\left(\dfrac{7}{4}\right)^2 + \left(\dfrac{7}{3}\right)^2} = 7\sqrt{\dfrac{3^2+4^2}{(3^2)(4^2)}} = 7\cdot \dfrac{5}{12} = \boxed{\dfrac{35}{12} \textbf{ (C)}}\]



	\item (\textit{AMC 12 12B Fall 2021}) What is the average number of pairs of consecutive integers in a randomly selected subset of $5$ distinct integers chosen from the set $\{ 1, 2, 3, \ldots, 30\}$? (For example the set $\{1, 17, 18, 19, 30\}$ has $2$ pairs of consecutive integers.)


 $\textbf{(A)}\ \frac{2}{3} \qquad\textbf{(B)}\ \frac{29}{36} \qquad\textbf{(C)}\ \frac{5}{6} \qquad\textbf{(D)}\ \frac{29}{30} \qquad\textbf{(E)}\ 1$

 \textbf{Answer: }A



\textbf{Solution 1}

There are $29$ possible pairs of consecutive integers, namely $p_1=\{1,2\},  \cdots, p_{29}=\{29,30\}$. Define a random variable $X_i$, with $X_i=1$, if $p_i$ is part of the 5-element subset, and $0$ otherwise. Then the number of pairs of consecutive integers in a $5$-element selection is given by the sum $X_1+\cdots + X_{29}$. By linearity of expectation, the expected value is equal to the sum of the $\mathbb{E}[X_i]$:
 \[\mathbb{E}[X_1+\cdots +X_{29}]=\mathbb{E}[X_1]+\cdots + \mathbb{E}[X_{29}]\]To compute $\mathbb{E}[X_i]$, note that $X_i=1$ for a total of $\binom{28}{3}$ out of $\binom{30}{5}$ possible selections. Thus
 \[\mathbb{E}[X_i]=\frac{\binom{28}{3}}{\binom{30}{5}}= \frac 1{29}\cdot \frac 23, \quad \textrm{which implies} \quad \mathbb{E}[X_1+\cdots +X_{29}]= \boxed{\textbf{(A)}\ \frac{2}{3}}.\]~kingofpineapplz~minor edit by sujoybec51

 



\textbf{Solution 2}

Alternatively, using linearity of expectation on the chosen subset: there are ${5 \choose 2}$ pairs of integers. There are 29 pairs of consecutive integers out of ${30 \choose 2}$ possible pairs, thus the probability that any given pair is consecutive is $\frac{29}{{30 \choose 2}}$. By linearity of expectation, there are $\frac{{5 \choose 2}29}{{30 \choose 2}} = \frac{10 \cdot 29}{\frac{30 \cdot 29}{2}} = \boxed{\textbf{(A)}\ \frac{2}{3}}$ such pairs on average.

 



\textbf{Solution 3 (casework bash)}

We will proceed with some casework. Let $n$ be the number of sets of consecutive numbers in the subset.Note that the maximum possible value of $n$ is $4.$Case $1$: $n = 4$There is only one way to arrange the distribution of the number of elements in each block here. There are $\dbinom{26}{1} = 26$ ways to arrange where that block of $5$ numbers will go, so a total of $1 \cdot 26 = 26$ ways for this case.Case $2$: $n = 3$There are $4$ ways to arrange the distribution of the number of elements in each block here. (4-1, 3-2, 2-3, 1-4). There are $\dbinom{26}{2}$ ways to arrange where those two blocks of numbers go among the other numbers such that no two of those blocks are adjacent to each other. Total for this case: $4 \cdot 325 = 1300.$Case $3$: $n = 2$By Stars and Bars, there are $\dbinom{4}{2} = 6$ ways to arrange the distribution of the number of elements in each block here. There are $\dbinom{26}{3}$ ways to arrange where those three blocks of numbers go among the other numbers such that no two of those blocks are adjacent to each other. Total for this case: $6 \cdot 2600 = 15600.$Case $4$: $n = 1$By Stars and Bars, there are $\dbinom{4}{3} = 4$ ways to arrange the distribution of the number of elements in each block here. There are $\dbinom{26}{4}$ ways to arrange where those four blocks of numbers go among the other numbers such that no two of those blocks are adjacent to each other. Total for this case: $4 \cdot 14950 = 59800.$Since the last case $n=0$ doesn't contribute to the expected value sum, we can ignore it. Using the expected value formula, our desired value is 
 \[\frac{4 \cdot 26 + 3 \cdot 1300 + 2 \cdot 15600 + 1 \cdot 59800}{\tbinom{30}{5}} = \frac{95004}{142506} = \boxed{\frac{2}{3}}.\]

 -fidgetboss\_4000

 



\textbf{Solution 4}

We define an outcome as $\left( a_1 ,\cdots, a_5 \right)$ with $1 \leq a_1 < a_2 < a_3 < a_4 < a_5 \leq 30$.

 We denote by $\Omega$ the sample space. Hence. $| \Omega | = \binom{30}{5}$.

 $\textbf{Case 1}$: There is only 1 pair of consecutive integers.

 $\textbf{Case 1.1}$: $\left( a_1 , a_2 \right)$ is the single pair of consecutive integers.

 We denote by $E_{11}$ the collection of outcomes satisfying this condition.Hence, $| E_{11} |$ is the number of outcomes satisfying
 \[ \left\{ \begin{array}{l} a_1 \geq 1 \\ a_3 \geq a_1 + 3 \\ a_4 \geq a_3 + 2 \\ a_5 \geq a_4 + 2 \\ a_5 \leq 30 \\ a_1, a_3, a_4, a_5 \in \Bbb N \end{array} \right.. \]

 Denote $b_1 = a_1 - 1$, $b_2 = a_3 - a_1 - 3$, $b_3 = a_4 - a_3 - 2$, $b_4 = a_5 - a_4 - 2$, $b_5 = 30 - a_5$.Hence,$| E_{11} |$ is the number of outcomes satisfying
 \[ \left\{ \begin{array}{l} b_1 + b_2 + b_3 + b_4 + b_5 = 22 \\ b_1, b_2 , b_3, b_4, b_5 \mbox{ are non-negative integers } \end{array} \right.. \]

 Therefore, $| E_{11} | = \binom{22 + 5 - 1}{5 - 1} = \binom{26}{4}$.

 $\textbf{Case 1.2}$: $\left( a_2 , a_3 \right)$ is the single pair of consecutive integers.

 We denote by $E_{12}$ the collection of outcomes satisfying this condition.Hence, $| E_{12} |$ is the number of outcomes satisfying
 \[ \left\{ \begin{array}{l} a_1 \geq 1 \\ a_2 \geq a_1 + 2 \\ a_4 \geq a_2 + 3 \\ a_5 \geq a_4 + 2 \\ a_5 \leq 30 \\ a_1, a_2, a_4, a_5 \in \Bbb N \end{array} \right.. \]

 Similar to our analysis for Case 1.1,  $| E_{12} | = \binom{22 + 5 - 1}{5 - 1} = \binom{26}{4}$.

 $\textbf{Case 1.3}$: $\left( a_3 , a_4 \right)$ is the single pair of consecutive integers.

 We denote by $E_{13}$ the collection of outcomes satisfying this condition.Hence, $| E_{13} |$ is the number of outcomes satisfying
 \[ \left\{ \begin{array}{l} a_1 \geq 1 \\ a_2 \geq a_1 + 2 \\ a_3 \geq a_2 + 2 \\ a_5 \geq a_3 + 3 \\ a_5 \leq 30 \\ a_1, a_2, a_3, a_5 \in \Bbb N \end{array} \right.. \]

 Similar to our analysis for Case 1.1,  $| E_{13} | = \binom{22 + 5 - 1}{5 - 1} = \binom{26}{4}$.

 $\textbf{Case 1.4}$: $\left( a_4 , a_5 \right)$ is the single pair of consecutive integers.

 We denote by $E_{14}$ the collection of outcomes satisfying this condition.Hence, $| E_{14} |$ is the number of outcomes satisfying
 \[ \left\{ \begin{array}{l} a_1 \geq 1 \\ a_2 \geq a_1 + 2 \\ a_3 \geq a_2 + 2 \\ a_4 \geq a_3 + 2 \\ a_4 \leq 29 \\ a_1, a_2, a_3, a_4 \in \Bbb N \end{array} \right.. \]

 Similar to our analysis for Case 1.1,  $| E_{14} | = \binom{22 + 5 - 1}{5 - 1} = \binom{26}{4}$.

 $\textbf{Case 2}$: There are 2 pairs of consecutive integers.

 $\textbf{Case 2.1}$: $\left( a_1 , a_2 \right)$ and $\left( a_2 , a_3 \right)$ are two pairs of consecutive integers.

 We denote by $E_{21}$ the collection of outcomes satisfying this condition.Hence, $| E_{21} |$ is the number of outcomes satisfying
 \[ \left\{ \begin{array}{l} a_1 \geq 1 \\ a_4 \geq a_1 + 4 \\ a_5 \geq a_4 + 2 \\ a_5 \leq 30 \\ a_1, a_4, a_5 \in \Bbb N \end{array} \right.. \]

 Similar to our analysis for Case 1.1,  $| E_{21} | = \binom{23 + 4 - 1}{4 - 1} = \binom{26}{3}$.

 $\textbf{Case 2.2}$: $\left( a_1 , a_2 \right)$ and $\left( a_3 , a_4 \right)$ are two pairs of consecutive integers.

 We denote by $E_{22}$ the collection of outcomes satisfying this condition.Hence, $| E_{22} |$ is the number of outcomes satisfying
 \[ \left\{ \begin{array}{l} a_1 \geq 1 \\ a_3 \geq a_1 + 3 \\ a_5 \geq a_3 + 3 \\ a_5 \leq 30 \\ a_1, a_3, a_5 \in \Bbb N \end{array} \right.. \]

 Similar to our analysis for Case 1.1,  $| E_{22} | = \binom{23 + 4 - 1}{4 - 1} = \binom{26}{3}$.

 $\textbf{Case 2.3}$: $\left( a_1 , a_2 \right)$ and $\left( a_4 , a_5 \right)$ are two pairs of consecutive integers.

 We denote by $E_{23}$ the collection of outcomes satisfying this condition.Hence, $| E_{23} |$ is the number of outcomes satisfying
 \[ \left\{ \begin{array}{l} a_1 \geq 1 \\ a_3 \geq a_1 + 3 \\ a_4 \geq a_3 + 2 \\ a_4 \leq 29 \\ a_1, a_3, a_4 \in \Bbb N \end{array} \right.. \]

 Similar to our analysis for Case 1.1,  $| E_{23} | = \binom{23 + 4 - 1}{4 - 1} = \binom{26}{3}$.

 $\textbf{Case 2.4}$: $\left( a_2 , a_3 \right)$ and $\left( a_3 , a_4 \right)$ are two pairs of consecutive integers.

 We denote by $E_{24}$ the collection of outcomes satisfying this condition.Hence, $| E_{24} |$ is the number of outcomes satisfying
 \[ \left\{ \begin{array}{l} a_1 \geq 1 \\ a_2 \geq a_1 + 2 \\ a_5 \geq a_2 + 4 \\ a_5 \leq 30 \\ a_1, a_2, a_5 \in \Bbb N \end{array} \right.. \]

 Similar to our analysis for Case 1.1,  $| E_{24} | = \binom{23 + 4 - 1}{4 - 1} = \binom{26}{3}$.

 $\textbf{Case 2.5}$: $\left( a_2 , a_3 \right)$ and $\left( a_4 , a_5 \right)$ are two pairs of consecutive integers.

 We denote by $E_{25}$ the collection of outcomes satisfying this condition.Hence, $| E_{25} |$ is the number of outcomes satisfying
 \[ \left\{ \begin{array}{l} a_1 \geq 1 \\ a_2 \geq a_1 + 2 \\ a_4 \geq a_2 + 3 \\ a_4 \leq 29 \\ a_1, a_2, a_4 \in \Bbb N \end{array} \right.. \]

 Similar to our analysis for Case 1.1,  $| E_{25} | = \binom{23 + 4 - 1}{4 - 1} = \binom{26}{3}$.

 $\textbf{Case 2.6}$: $\left( a_3 , a_4 \right)$ and $\left( a_4 , a_5 \right)$ are two pairs of consecutive integers.

 We denote by $E_{26}$ the collection of outcomes satisfying this condition.Hence, $| E_{26} |$ is the number of outcomes satisfying
 \[ \left\{ \begin{array}{l} a_1 \geq 1 \\ a_2 \geq a_1 + 2 \\ a_3 \geq a_2 + 2 \\ a_3 \leq 28 \\ a_1, a_2, a_3 \in \Bbb N \end{array} \right.. \]

 Similar to our analysis for Case 1.1,  $| E_{26} | = \binom{23 + 4 - 1}{4 - 1} = \binom{26}{3}$.

 $\textbf{Case 3}$: There are 3 pairs of consecutive integers.

 $\textbf{Case 3.1}$: $\left( a_1 , a_2 \right)$, $\left( a_2 , a_3 \right)$ and $\left( a_3 , a_4 \right)$ are three pairs of consecutive integers.

 We denote by $E_{31}$ the collection of outcomes satisfying this condition.Hence, $| E_{31} |$ is the number of outcomes satisfying
 \[ \left\{ \begin{array}{l} a_1 \geq 1 \\ a_5 \geq a_1 + 5 \\ a_5 \leq 30 \\ a_1, a_5 \in \Bbb N \end{array} \right.. \]

 Similar to our analysis for Case 1.1,  $| E_{31} | = \binom{24 + 3 - 1}{3 - 1} = \binom{26}{2}$.

 $\textbf{Case 3.2}$: $\left( a_1 , a_2 \right)$, $\left( a_2 , a_3 \right)$ and $\left( a_4 , a_5 \right)$ are three pairs of consecutive integers.

 We denote by $E_{32}$ the collection of outcomes satisfying this condition.Hence, $| E_{32} |$ is the number of outcomes satisfying
 \[ \left\{ \begin{array}{l} a_1 \geq 1 \\ a_4 \geq a_1 + 4 \\ a_4 \leq 29 \\ a_1, a_4 \in \Bbb N \end{array} \right.. \]

 Similar to our analysis for Case 1.1,  $| E_{32} | = \binom{24 + 3 - 1}{3 - 1} = \binom{26}{2}$.

 $\textbf{Case 3.3}$: $\left( a_1 , a_2 \right)$, $\left( a_3 , a_4 \right)$ and $\left( a_4 , a_5 \right)$ are three pairs of consecutive integers.

 We denote by $E_{33}$ the collection of outcomes satisfying this condition.Hence, $| E_{33} |$ is the number of outcomes satisfying
 \[ \left\{ \begin{array}{l} a_1 \geq 1 \\ a_3 \geq a_1 + 3 \\ a_3 \leq 28 \\ a_1, a_3 \in \Bbb N \end{array} \right.. \]

 Similar to our analysis for Case 1.1,  $| E_{33} | = \binom{24 + 3 - 1}{3 - 1} = \binom{26}{2}$.

 $\textbf{Case 3.4}$: $\left( a_2 , a_3 \right)$, $\left( a_3 , a_4 \right)$ and $\left( a_4 , a_5 \right)$ are three pairs of consecutive integers.

 We denote by $E_{34}$ the collection of outcomes satisfying this condition.Hence, $| E_{34} |$ is the number of outcomes satisfying
 \[ \left\{ \begin{array}{l} a_1 \geq 1 \\ a_2 \geq a_1 + 2 \\ a_2 \leq 27 \\ a_1, a_2 \in \Bbb N \end{array} \right.. \]

 Similar to our analysis for Case 1.1,  $| E_{34} | = \binom{24 + 3 - 1}{3 - 1} = \binom{26}{2}$.

 $\textbf{Case 4}$: There are 4 pairs of consecutive integers.

 In this case, $\left( a_1, a_2 , a_3 , a_4 , a_5 \right)$ are consecutive integers.

 We denote by $E_4$ the collection of outcomes satisfying this condition.Hence, $| E_4 |$ is the number of outcomes satisfying
 \[ \left\{ \begin{array}{l} a_1 \geq 1 \\ a_1 \leq 27 \\ a_1 \in \Bbb N \end{array} \right.. \]

 Hence, $| E_4 | = 26$.

 Therefore, the average number of pairs of consecutive integers is
 \begin{align*} & \frac{1}{| \Omega|} \left( 1 \cdot \sum_{i=1}^4 | E_{1i} | + 2 \cdot \sum_{i=1}^6 | E_{2i} | + 3 \cdot \sum_{i=1}^4 | E_{3i} | + 4 \cdot  | E_4 | \right) \\ & = \frac{1}{\binom{30}{5}} \left( 4 \binom{26}{4} + 12 \binom{26}{3} + 12 \binom{26}{2} + 4 \cdot 26 \right) \\ & = \frac{2}{3} . \end{align*}


 Therefore, the answer is $\boxed{\textbf{(A) }\frac{2}{3}}$.

 



\textbf{Solution 5}

Let $a_1, a_2, a_3, a_4, a_5$ be the five numbers chosen. Then, we can consider the first number and the four differences, $a_1, a_2-a_1, a_3-a_2, a_4-a_3, a_5-a_4$. Now, each must be positive integers and the sum is less than or equal to $30$, which there are $\binom{30}{5}$ by Hockey-Stick Identity. Now, we can use the Linearity of Expectation on $a_2-a_1, a_3-a_2, a_4-a_3, a_5-a_4$ since each of them being $1$ represents a pair of consecutive integer. For each of them, we have $\binom{28}{3}+\binom{27}{3}+...+1=\binom{29}{4}$. Now, we have four such numbers to consider, and by Linearity of Expectation, its $4\cdot \binom{29}{4}$. Now, we divide by $\binom{30}{5}$ to get the answer of $\boxed{\textbf{(A)} ~\frac{2}{3}}$.

 



\textbf{Solution 6}

We will first list out the possible pairs: 
 \[(1,2), (2,3), (3,4), (4,5), \ldots (29,30)\]. Now, we can find the probability of each one and multiply it by $29$ because of Linearity of Expectation. So, we have: 
 \[29 \cdot \dfrac{\binom{28}{3}}{\binom{30}{5}} = \boxed{\textbf{(A)} ~\frac{2}{3}}\]



	\item (\textit{AMC 12 12B Fall 2021}) Triangle $ABC$ has side lengths $AB = 11, BC=24$, and $CA = 20$. The bisector of $\angle{BAC}$ intersects $\overline{BC}$ in point $D$, and intersects the circumcircle of $\triangle{ABC}$ in point $E \ne A$. The circumcircle of $\triangle{BED}$ intersects the line $AB$ in points $B$ and $F \ne B$. What is $CF$?


 $\textbf{(A) } 28 \qquad \textbf{(B) } 20\sqrt{2} \qquad \textbf{(C) } 30 \qquad \textbf{(D) } 32 \qquad \textbf{(E) } 20\sqrt{3}$

 \textbf{Answer: }C



\textbf{Solution 1 (Reflection)}


                    \begin{figure}[H]
                        \centering
                        \includegraphics[width=0.5\textwidth]{Solution/AMC_12_12B_Fall_2021_24_Solution_Figure_1.png}
                    \end{figure}
                    By the Law of Cosine $\cos A = \frac{AC^2 + AB^2 - BC^2}{ 2 \cdot AC \cdot AB} = \frac{20^2 + 11^2 - 24^2}{2\cdot20\cdot11} = -\frac18$

 As $ABEC$ is a cyclic quadrilateral, $\angle CEA = \angle CBA$. As $BDEF$ is a cyclic quadrilateral, $\angle CBA = \angle FEA$.

 $\because \quad \angle CEA = \angle FEA \quad \text{and} \quad \angle CAE = \angle FAE$
$\therefore \quad \triangle AFE \cong \triangle ACE$ by $ASA$

 Hence, $AF = AC = 20$

 By the Law of Cosine $CF = \sqrt{20^2 + 20^2 - 2 \cdot 20 \cdot 20 (-\frac18)} = \sqrt{900} = \boxed{\textbf{C}~\text{30}}$

 Note that $F$ is $C$'s reflection over line $AE$, quadrilateral $ACEF$ is a kite symmetrical by line $AE$, $AE \perp CF$.

 



\textbf{Solution 2 (Olympiad Solution using Spiral Similarity)}

Construct the $E$-antipode, $E^{\prime}\in(ABC)$. Notice $\triangle CE^{\prime}A\stackrel{+}{\sim}\triangle CBF$ by spiral similarity at $C$, thus $CF=\dfrac{CB\cdot CA}{CE^{\prime}}=\frac{480}{CE^{\prime}}$. Let $CE^{\prime}=x$; by symmetry $BE^{\prime}=x$ as well and $\cos\angle BE^{\prime}C=\cos\angle A=\frac{11^{2}+20^{2}-24^{2}}{2\cdot 11\cdot 20}=-\frac{1}{8}$ from Law of Cosines in $\triangle ABC$, so by Law of Cosines in $\triangle BE^{\prime}C$ we have 
 \[x^{2}+x^{2}+\left(2x^{2}\right)\left(-\dfrac{1}{8}\right)=24^{2}\] from which $x=16$. Now, $CF=\dfrac{480}{16}=\boxed{\textbf{C}~\text{30}}$.

 
                    \begin{figure}[H]
                        \centering
                        \includegraphics[width=0.5\textwidth]{Solution/AMC_12_12B_Fall_2021_24_Solution_Figure_2.png}
                    \end{figure}
                    




\textbf{Solution 3}

\textbf{Claim:} $\triangle ADC \sim \triangle ABE.$

 \textbf{Proof:} Note that $\angle CAD = \angle CAE = \angle EAB$ and $\angle DCA = \angle BCA = \angle BEA$ meaning that our claim is true by AA similarity.

 Because of this similarity, we have that 
 \[\frac{AC}{AD} = \frac{AE}{AB} \Longrightarrow AB \cdot AC = AD \cdot AE = AB \cdot AF\] by Power of a Point. Thus, $AC=AF=20.$

 Two solution methods follow from here.

 



\textbf{Solution 3.1 (Stewart's theorem)}

Applying \href{https://artofproblemsolving.com/wiki/index.php?title=Stewart\%27s\_theorem}{Stewart's theorem} on $\triangle ABC$ with cevian $\overline{CF}$ using the \href{https://artofproblemsolving.com/wiki/index.php?title=Directed\_legnths}{directed lengths} $AF = AC = 20$ and $FB = 11-20 = -9$, we obtain 
 \begin{align*} (20)(-9)(11) + (CF)(11)(CF) &= (24)(20)(24) + (20)(-9)(20) \\ 11CF^{2} - 1980 &= 11520 - 3600\end{align*}
 so $CF=\sqrt{\frac{11520 - 3600 + 1980}{11}}=\sqrt{\frac{9900}{11}}=\sqrt{900}=\boxed{\textbf{(C) }30}$.

 



\textbf{Solution 3.2 (Double Cosine Law)}

Note that $\angle CAF = \angle CAB$ so we may plug into Law of Cosines to find the angle's cosine: 
 \[AB^2+AC^2-2\cdot AB \cdot AC \cdot \cos(\angle CAB) = BC^2 \to \cos(\angle CAB) = -\frac{1}{8}.\]

 So, we observe that we can use Law of Cosines again to find $CF$: 
 \[CF^2 = AF^2+AC^2-2 \cdot AF \cdot AC \cdot \cos(\angle CAF) = 900 \to CF=\boxed{\textbf{(C) }30}\] both ways.

 - Kevinmathz

 



\textbf{Solution 4}

This solution is based on this figure: 
                    \begin{figure}[H]
                        \centering
                        \includegraphics[width=0.5\textwidth]{Solution/AMC_12_12B_Fall_2021_24_Solution_Figure_3.png}
                    \end{figure}
                    

 Denote by $O$ the circumcenter of $\triangle BED$.Denote by $R$ the circumradius of $\triangle BED$.

 In $\triangle BCF$, following from the law of cosines, we have
 \begin{align*} CF^2 & = BC^2 + BF^2 - 2 BC \cdot BF \cos \angle CBF \\ & = BC^2 + BF^2 + 2 BC \cdot BF \cos \angle ABC .  \hspace{1cm} (1) \end{align*}
For $BF$, we have
 \begin{align*} BF & = 2 R \cos \angle FBO \\ & = 2 R \cos \left( 180^\circ - \angle ABC - \angle CBO \right) \\ & = 2 R \cos \left( 180^\circ - \angle ABC - \frac{180^\circ - \angle BOD}{2} \right) \\ & = 2 R \cos \left( 180^\circ - \angle ABC - \frac{180^\circ - 2 \angle BED}{2} \right) \\ & = 2 R \cos \left( 180^\circ - \angle ABC - \frac{180^\circ - 2 \angle BCA}{2} \right) \\ & = 2 R \cos \left( 90^\circ - \angle ABC + \angle BCA \right) \\ & = 2 R \sin \left( \angle ABC - \angle BCA \right) \\ & = \frac{BD}{\sin \angle BED} \sin \left( \angle ABC - \angle BCA \right) \\ & = \frac{BD}{\sin \angle BCA} \sin \left( \angle ABC - \angle BCA \right) \\ & = BD \left( \sin \angle ABC \cot \angle BCA - \cos \angle ABC \right) . \hspace{1cm} (2) \end{align*}
The fourth equality follows from the property that $B$, $D$, $E$ are concyclic.The fifth and the ninth equalities follow from the property that $A$, $B$, $C$, $E$ are concyclic.

 Because $AD$ bisects $\angle BAC$, following from the angle bisector theorem, we have
 \[ \frac{BD}{CD} = \frac{AB}{AC} . \]Hence, $BD = \frac{24 \cdot 11}{31}$.

 In $\triangle ABC$, following from the law of cosines, we have
 \begin{align*} \cos \angle ABC & = \frac{AB^2 + BC^2 - AC^2}{2 AB \cdot BC} \\ & = \frac{9}{16} \end{align*}
and
 \begin{align*} \cos \angle BCA & = \frac{AC^2 + BC^2 - AB^2}{2 AC \cdot BC} \\ & = \frac{57}{64} . \end{align*}
Hence, $\sin \angle ABC = \frac{5 \sqrt{7}}{16}$ and $\sin \angle BCA = \frac{11 \sqrt{7}}{64}$.Hence, $\cot \angle BCA = \frac{57}{11 \sqrt{7}}$.

 Now, we are ready to compute $BF$ whose expression is given in Equation (2).We get $BF = 9$.

 Now, we can compute $CF$ whose expression is given in Equation (1).We have $CF = 30$.

 Therefore, the answer is $\boxed{\textbf{(C) }30}$.

 



\textbf{Solution 5}

Denote $B=(0, 0)$ and $C=(24, 0)$. Note that by Heron's formula the area of $\triangle ABC$ is $\frac{165\sqrt{7}}{4}$ so the $y$-coordinate of $A$ (height of $A$ above the $x$-axis) is easily computed by the base-height formula as $\frac{55\sqrt7}{16}$.

 Now, since $AB=11$, the $x$-coordinate of $A$ satisfies $x^2+(\frac{55\sqrt7}{16})^2=11^2$ and solving gives $x=\frac{99}{16}$.

 The circumcircle of $\triangle ABC$ has radius $\frac{abc}{4A}=\frac{11\cdot 24\cdot 20}{165\sqrt7}=\frac{32}{\sqrt7}$. We know by the perpendicular bisector rule that the circumcenter $O$ is located directly below the midpoint of $\overline{BC}$ ($x$-coordinate $12$).

 So, the negative $y$ coordinate of $O$ satisfies $12^2+y^2=(\frac{32}{\sqrt7})^2$ and solving gives $y=-\frac{4}{\sqrt7}$.

 It's also clear that point $E$ is going to be located directly below $O$ on the circle, because the angle bisector intersects the circumcircle at the midpoint of the arc (Fact 5). Since the radius of the circle is $\frac{32}{\sqrt7}$, we have the coordinates of $E=(12, -\frac{36}{\sqrt7})$

 Solving for point $D$ (the point on the $x$-axis between $A$ and $E$), we get that $D=(\frac{264}{31}, 0)$.

 So now we know six of the critical points: $A=(\frac{99}{16}, \frac{55\sqrt7}{16})$; $B=(0, 0)$; $C=(24, 0)$; $D=(\frac{264}{31}, 0)$; $E=(12, -\frac{36}{\sqrt7})$; $O=(12, -\frac{4}{\sqrt7})$.

 We are now ready to add in the circumcircle of $\triangle BDE$, which has radius $\frac{BD\cdot DE\cdot BE}{4[BDE]}$. From the above information, $BD=\frac{264}{31}$, $DE=\sqrt{(\frac{108}{31})^2+(\frac{36}{\sqrt7})^2}$, and $BE=\sqrt{12^2+(\frac{36}{\sqrt7})^2}$.

 After a bit of simplification we end up with $DE=\frac{1152}{31\sqrt7}$ and $BE=\frac{48}{\sqrt7}$.

 For the area of $\triangle BDE$, the altitude dropped from vertex $E$ has height $\frac{36}{\sqrt7}$, and the base $\overline{BD}$ has length $\frac{264}{31}$, so its area is $\frac12\cdot\frac{36}{\sqrt7}\cdot\frac{264}{31}=\frac{4752}{31\sqrt7}$.

 Thus, $\frac{BD\cdot DE\cdot BE}{4[BDE]}=\frac{\frac{264}{31}\cdot\frac{1152}{31\sqrt7}\cdot\frac{48}{\sqrt{7}}}{4\cdot\frac{4752}{31\sqrt7}}$ which after tons of cancellations becomes $\frac{768}{31\sqrt7}$.

 We know from the perpendicular bisector rule that the circumcenter $P$ of $\triangle BDE$ is located directly below the midpoint of $\overline{BD}$ ($x$-coordinate $\frac{132}{31}$).

 So, the negative $y$-coordinate of $P$ satisfies $(\frac{132}{31})^2+y^2=(\frac{768}{31\sqrt7})^2$, and solving gives $y=-\frac{684}{31\sqrt7}$. Thus, the equation of the circumcircle of $\triangle BDE$ is $(x-\frac{132}{31})^2+(y+\frac{684}{31\sqrt7})^2=(\frac{768}{31\sqrt7})^2$.

 Point $F$ is the intersection of this circle and the line $\overline{AB}$, which has equation $y=\frac{5\sqrt7}{9}x$. So, we substitute $y=\frac{5\sqrt7}{9}x$ into the equation of the circle to get $(x-\frac{132}{31})^2+(\frac{5\sqrt7}{9}x+\frac{684}{31\sqrt7})^2=(\frac{768}{31\sqrt7})^2$.

 After simplifying, we have $\frac{256}{81}x^2+16x=0$ (the $\frac{768}{31\sqrt7}$'s cancel out), whose solutions are $x=0$ and $x=-\frac{81}{16}$. The first corresponds to the origin, and the second corresponds to point $F$. Thus the coordinates of $F$ are $(-\frac{81}{16}, \frac{5\sqrt7}{9}\cdot\frac{-81}{16})=(-\frac{81}{16}, -\frac{45\sqrt7}{16})$.

 The coordinates of $C$ are $(24, 0)$, so 
 \[CF=\sqrt{(24+\frac{81}{16})^2+(\frac{45\sqrt7}{16})^2}=\sqrt{(\frac{465}{16})^2+(\frac{45\sqrt7}{16})^2}=\frac{\sqrt{465^2+(45\sqrt7)^2}}{16}=\frac{\sqrt{(15\cdot 31)^2+(15\cdot 3\sqrt7)^2}}{16}=\frac{15\sqrt{31^2+(3\sqrt7)^2}}{16}=\frac{15\sqrt{961+63}}{16}=\frac{15\sqrt{1024}}{16}=\frac{15}{16}\cdot 32=30.\]



	\item (\textit{AMC 12 12B Fall 2021}) For $n$ a positive integer, let $R(n)$ be the sum of the remainders when $n$ is divided by $2$, $3$, $4$, $5$, $6$, $7$, $8$, $9$, and $10$. For example, $R(15) = 1+0+3+0+3+1+7+6+5=26$. How many two-digit positive integers $n$ satisfy $R(n) = R(n+1)\,?$
$\textbf{(A) }0\qquad\textbf{(B) }1\qquad\textbf{(C) }2\qquad\textbf{(D) }3\qquad\textbf{(E) }4$

 \textbf{Answer: }C



\textbf{Solution 1}

Note that we can add $9$ to $R(n)$ to get $R(n+1)$, but must subtract $k$ for all $k|n+1$. Hence, we see that there are four ways to do that because $9=7+2=6+3=5+4=4+3+2$. Note that only $7+2$ is a plausible option, since $4+3+2$ indicates $n+1$ i

 


\end{enumerate}